% ====================================================================
% graphite_ica_ch1_rebuilt.tex — Chapter 1 (RB 재구성본)
%   흑연 음극 ICA(dQ/dV) tail: 열역학 무대 + 극판 전위에 의한 배리어 낮춤
%   → 진행률 완화 → relaxation-length spectrum → 피팅 가능한 닫힌 논리식.
%
%   재구성 원칙(RB plan 2026-05-30):
%     · 학부생 가독(GS-1,2): 모든 식은 전제부터 무비약 유도.
%     · 수식 흐름 보존(방향-2): consolidated_ch1 의 식 순서·역할 유지.
%     · 모든 가정 = 통합 Assumption Ledger(AL-#) + 논문/교과서 근거(GS-3).
%     · 계산편의 비약 배제(GS-4): 임의 cap/clip/step/softplus/근거없는 임계값 제거.
%     · closure(Plan A/B)·수치·피팅 실무를 본 장에 통합(Ch6 해체, 2026-06-01): §planB/§planA/§ch1_simplefit/부록.
%   AL 번호 = 통합 master(RB_AL_MASTER.md) global 체계.
% Date: 2026-05-30  Author: Project_Anode_Fit (Claude 측, RB-series)
% ====================================================================
\documentclass[11pt,a4paper]{article}
\usepackage[margin=22mm]{geometry}
\usepackage{kotex}
\usepackage{amsmath,amssymb,mathtools,bm}
\usepackage{booktabs,longtable,array}
\usepackage{enumitem}
\usepackage{amsthm}
\usepackage{hyperref}
\usepackage{fancyhdr}
\usepackage{lastpage}
\usepackage{setspace}
\setstretch{1.13}
\setlength{\parskip}{0.45em}\setlength{\parindent}{0pt}\setlength{\headheight}{14pt}
\hypersetup{colorlinks=true, linkcolor=blue, urlcolor=blue, citecolor=blue,
  pdftitle={흑연 음극 ICA tail — Chapter 1 (RB 재구성본)}}
\pagestyle{fancy}\fancyhf{}
\lhead{흑연 음극 ICA tail — Ch.1 (열역학 무대 + 배리어 낮춤)}\rhead{\thepage/\pageref{LastPage}}
\renewcommand{\headrulewidth}{0.4pt}

\newtheorem*{groundingbox}{Grounding}
\newtheorem*{boundbox}{유효범위(Validity bound)}
\newtheorem*{flagbox}{FLAGGED (미확정)}
\newtheorem*{keybox}{Keystone}
\newtheorem*{linkbox}{Ch1--5 전달식}
\newtheorem*{verifybox}{유도 검증 (차원·극한·부호)}
\newtheorem*{practicebox}{실무 팁 (본문 물리식 아님)}
\newtheorem*{loopbox}{도입 관찰로의 폐합}

\newcommand{\dd}{\mathrm{d}}
\newcommand{\eff}{\mathrm{eff}}
\newcommand{\eq}{\mathrm{eq}}
\newcommand{\app}{\mathrm{app}}
\newcommand{\drive}{\mathrm{drive}}
\newcommand{\bg}{\mathrm{bg}}
\newcommand{\cell}{\mathrm{cell}}
\newcommand{\ext}{\mathrm{ext}}
\newcommand{\OCV}{\mathrm{OCV}}
\newcommand{\tot}{\mathrm{tot}}
\newcommand{\irr}{\mathrm{irr}}
\newcommand{\rev}{\mathrm{rev}}
\newcommand{\ct}{\mathrm{ct}}
\newcommand{\transp}{\mathrm{transport}}
\newcommand{\cb}{\mathrm{cb}}
\newcommand{\dyn}{\mathrm{dyn}}
\newcommand{\sstat}{\mathrm{ss}}
\newcommand{\conv}{\mathrm{conv}}
\newcommand{\stt}{\mathrm{state}}
\newcommand{\hys}{\mathrm{hys}}
\newcommand{\tar}{\mathrm{tar}}
\newcommand{\loss}{\mathrm{loss}}
\newcommand{\obs}{\mathrm{obs}}
\newcommand{\tol}{\mathrm{tol}}
\newcommand{\noise}{\mathrm{noise}}
\newcommand{\brk}{\discretionary{}{}{}}% 표/참고문헌 긴 DOI 줄바꿈 허용(zero-width)
\newcommand{\AL}[1]{\textsf{[AL-#1]}}

\title{\textbf{리튬이온전지 흑연 음극 ICA($dQ/dV$) tail 이론 — Chapter 1}\\[0.4em]
\large 열역학 무대($V_n$, 평형 target) + 극판 전위에 의한 배리어 낮춤(동역학 주역)\\[0.2em]
\normalsize 전하보존 $\to$ 유효배리어 $\to$ relaxation-length spectrum $\to$ 피팅 가능한 닫힌 논리식}
\author{Project\_Anode\_Fit (RB 재구성본)}
\date{2026-05-30}

\begin{document}
\maketitle

% ====================================================================
\section*{서 — 하나의 관찰에서 출발하는 단일 인과 사슬}
\addcontentsline{toc}{section}{서 — 단일 인과 사슬}
% ====================================================================
본 장은 다음 \textbf{하나의 관찰}에서 출발하여 끝까지 그 끈을 놓지 않는다.
\begin{quote}\itshape
흑연 음극의 $dQ/dV$(ICA, incremental capacity) 곡선에서 상변이 peak 의 꼬리(tail)가 \textbf{온도가
낮을수록 길게 늘어져 다음 peak 와 겹치고}, 높을수록 짧게 끝난다. 또 전류(C-rate)가 커지면 peak 겹침이
심해진다.
\end{quote}
이 관찰을 \textbf{다섯 단계}로 푼다. 각 단계는 학부 수준(화학열역학·전기화학·미적분)으로 따라갈 수 있게
중간을 생략하지 않는다.
\begin{enumerate}[label=\textbf{(\arabic*)},leftmargin=2.2em]
\item \textbf{평형 열역학만으로는 설명이 안 된다.} 뒤에서 보이듯 평형 전이폭은 $w=RT/F$ 라 \emph{저온서
  오히려 좁아진다}(\S\ref{sec:eq}). 평형 이론은 ``저온 $\to$ 짧은 꼬리''를 예측한다 — 관찰과 \emph{반대}.
\item \textbf{따라서 꼬리의 주역은 동역학이다.} 진행률 $\xi_j$ 가 평형 목표 $\xi_{\eq,j}$ 를 \emph{유한
  속도}로 좇으며 생기는 지연(lag)이 꼬리를 만든다. 속도상수 $k\propto e^{-\Delta G_a/RT}$ 가 저온서
  작아져 지연이 커지고 꼬리가 길어진다(\S\ref{sec:dyn}) — 관찰과 \emph{일치}.
\item \textbf{그 배리어는 극판 전위가 낮춘다.} 전류가 흐를 때 내부 전위는 평형(OCV)과 다르며, 그 전위
  구동력 $\mathcal A_j$ 가 유효 배리어를 낮춰 진행 속도를 바꾼다(\S\ref{sec:rate}). 그래서 C-rate 가
  peak 겹침을 키운다.
\item \textbf{무대는 전하 보존이 깐다.} 내부 음극 전위 $V_n$ 은 외부에서 읽어오는 값이 아니라 \emph{전하
  보존식의 해}로 결정된다(\S\ref{sec:charge}). 꼬리의 \emph{깊이}는 배리어 분포가 만드는
  \emph{relaxation-length spectrum} 의 중첩(kernel integral)으로 설명된다(\S\ref{sec:spectrum}--\ref{sec:kernel}).
\item \textbf{결과 = 피팅 가능한 닫힌 논리식}(\S\ref{sec:fiteq}). 이후 발열(Ch.2,4)$\cdot$반응속도론
  (Ch.3)$\cdot$히스테리시스(Ch.5)로 확장하고, 수치$\cdot$피팅 실무는 \emph{본 장 안}(\S\ref{sec:ch1_simplefit} 심플
근사식, \S\ref{sec:planB}--\ref{sec:planA} closure, 부록 \S\ref{sec:ch1_numeric})에서 닫는다.
\end{enumerate}
\noindent\textbf{한 문장 요지: 열역학이 \emph{무대}($V_n$ 과 평형 target)를 깔고, 관찰된 꼬리 거동의
\emph{주역}은 동역학이다.} 본 Chapter 1 은 \textbf{방전(탈리튬화) 방향}만 다룬다. 충전 branch,
히스테리시스, full-cell 전압 분해는 Ch.5 에서 정의한다.

\begin{groundingbox}
\textbf{지배 표준(GS).} \textbf{(GS-1)} 모든 식은 물리적 의미를 먼저 말로 제시한 뒤 수식화하고 중간
단계를 생략하지 않는다(학부 무비약). \textbf{(GS-2)} 결론은 출판 수준 엄밀성. \textbf{(GS-3)} 모든 가정은
통합 Assumption Ledger(\AL{\#}: 논문/교과서 근거 + 검증 DOI)를 인용하며, 근거가 없으면 FLAGGED 로
표기하고 established 로 쓰지 않는다. \textbf{(GS-4)} 임의 clip/cap/softplus/step-function 정의식과 근거
없는 수치 임계값을 물리 모델로 쓰지 않는다(속도 제한은 물리적 병목으로 도입). 수치 해법과 피팅 실무는
본 장 부록 \S\ref{sec:ch1_numeric} 로 분리한다.
\end{groundingbox}

\tableofcontents
\newpage

% ====================================================================
\section{기호와 단위}\label{sec:notation}
% ====================================================================
진행률 기호는 $\xi$ 로 통일한다(일부 문헌의 $\theta$ 와 동일물). 방전 누적 전하 좌표 $q$, 외부 전류 $I$,
방향 부호 $s_I,s_{\phi,j}\in\{+1,-1\}$(방전 기본 $+1$).
\begin{longtable}{p{0.16\textwidth} p{0.10\textwidth} p{0.66\textwidth}}
\toprule 기호 & 단위 & 의미 \\ \midrule \endhead
$q$ & --- & 방전 진행 좌표, $q=Q_\ext/Q_\cell$ \\
$Q_\ext,\ Q_\cell$ & C & 외부 누적 방전 전하, 기준 셀 용량(쿨롱) \\
$V_n$ & V & 내부 음극 전위 — 전하 보존식의 해(\AL{9}) \\
$V_{n,\app}$ & V & 측정되는 apparent 전위 $=V_n+s_I|I|R_n$ (\AL{6}) \\
$V_{n,\drive}$ & V & 속도식 구동 전위; 기본 $=V_n$ \\
$V_{n,\OCV}$ & V & 평형($\xi_j=\xi_{\eq,j}$)에서의 전하보존식 해 \\
$R_n$ & $\Omega$ & 음극 유효 분극 계수(전압축 이동) \\
$\xi_j,\ \xi_{\eq,j}$ & --- & $j$번째 effective transition 진행률, 평형 진행률 ($0\!\le\!\xi\!\le\!1$) \\
$U_j(T),\ w_j(T)$ & V & 전이 중심 전위(\S\ref{sec:eq} 에서 $U_j\equiv U_j^0-\mu^0/(s_{\phi,j}F)$ 로 $\mu^0$ 흡수), 평형 전이 폭 \\
$Q_{j,\tot}$ & C & $j$번째 transition 의 방전 방향 용량 기여 ($>0$) \\
$Q_\bg(V_n,T)$ & C & staging 으로 분리 안 되는 배경 누적 용량; $\partial Q_\bg/\partial V_n$=잔류 chemical capacitance \\
$\mathcal A_j$ & J/mol & 상변이 구동력 $=s_{\phi,j}F(V_{n,\drive}-U_j)$ \\
$\chi_j$ & --- & 배리어 낮춤 민감도 ($0\le\chi_j\le1$; Ch.3 transfer coefficient $\beta_j$ 와는 \emph{대칭 Marcus 1차$\cdot$정상 target 극한 한정}으로만 동일물 — 일반은 별개, \S\ref{sec:rate}) \\
$\Delta G_{a,j},\Delta H_{a,j},\Delta S_{a,j}$ & J/mol & 활성화 자유에너지/엔탈피/엔트로피 \\
$G$ & J/mol & local activation barrier (mode 별 값; \S\ref{sec:spectrum} 에서 분포) \\
$k_j,\ k_0(T)$ & 1/s & 진행률 완화 속도상수(mobility), Eyring prefactor \\
$L$ & --- & local relaxation length $=|I|/(Q_\cell k_j)$ \\
$\rho_G(G;T)$ & mol/J & activation barrier 확률밀도 ($\int_0^\infty\rho_G\,\dd G=1$; $G$ 가 J/mol 이므로 단위는 그 역수 mol/J) \\
$A_L^{\mathrm{prob}},A_L^{\mathrm{amp}}$ & --- & relaxation-length 확률밀도(정규화) / 진폭 스펙트럼 $=A_0 A_L^{\mathrm{prob}}$; kernel 의 $A_L\equiv A_L^{\mathrm{amp}}$ (\S\ref{sec:spectrum}) \\
$A_0(L)$ & --- & mode 별 물리적 진폭 가중(입자크기분포$\cdot$활성 site 밀도 등 독립 물리량) \\
$L_0,\ L_{\min}$ & --- & $L_0=|I|/(Q_\cell k_0)$ (prefactor 기준 길이), $L_{\min}=L_0 e^{-\chi_j\mathcal A_j/RT}$ (support 하한, $G\!\ge\!0$) \\
$L_\varphi$ & V & 전압축 특성 꼬리 길이 $=|\dd V_n/\dd q|_{q_a}L$ (\S\ref{sec:fiteq}, 식~\eqref{eq:Lphi}) \\
$\Theta,\ Q_p$ & ---,\ C & 전체 effective phase progress, 그 용량 scale \\
\bottomrule
\end{longtable}
\textbf{단위 주의.} 미분방정식의 전류 $I$ 가 A $=\mathrm{C/s}$ 이므로 $Q_\cell$ 은 쿨롱이어야 한다
($Q_\cell[\mathrm C]=3600\,Q_\cell[\mathrm{Ah}]$). 외부 전하·진행 좌표는 $Q_\ext(t)=\int_0^t|I|\,\dd t'$,
$q=Q_\ext/Q_\cell$, 정전류에서 $\dd q/\dd t=|I|/Q_\cell$.

% ====================================================================
\section{흑연 staging 과 effective transition}\label{sec:staging}
% ====================================================================
흑연의 lithiation 은 stage $4\to3\to2\mathrm L\to2\to1$ 의 순서로 진행하며, 각 전이가 $dQ/dV_n$ 에 peak 를
남긴다 \cite{dahn1991,ohzuku1993,funabiki1999ea,funabiki1999jes}(\AL{3}). 인접 전이의 중심 전위 차가
전이폭 이하이면 peak 가 융합되어 관측 peak 수가 $N_p\approx3\sim4$ 가 된다. 본 장은 분리 가능한 유효
transition 을 $j=1,\dots,N_p$ 로 둔다. 각 유효 transition 내부에도 입자 크기$\cdot$국소 환경$\cdot$결정립계
차이에 의한 \emph{local mode 분포}가 존재하며, 이것이 \S\ref{sec:spectrum} spectrum 의 물리적 근원이다.
진행률 방향은 방전 좌표 $q$ 와 같게 잡아 $\xi_j=0$(미진행)$\to\xi_j=1$(완료), $Q_{j,\tot}>0$.

% ====================================================================
\section{무대 (I): 전하 보존식과 내부 전위 $V_n$}\label{sec:charge}
% ====================================================================
\textbf{물리.} 외부 회로로 흘린 방전 전하와, 음극 내부의 방전 방향 누적 전하(배경 저장분 + staging 진행분)는
같아야 한다 — 단순한 전하 보존이다. 이 등식이 내부 전위 $V_n$ 을 결정한다(외부에서 주어지는 함수가
아니다; \AL{1},\AL{9}, \cite{doyle1993,newman}):
\begin{equation}
\boxed{\;Q_\cell\,q \;=\; Q_\bg(V_n,T) \;+\; \sum_{j=1}^{N_p} Q_{j,\tot}\,\xi_j\;.}
\label{eq:charge_balance}
\end{equation}
좌변은 외부 입력(아는 값), 우변은 내부 상태($V_n$ 과 $\xi_j$ 의 함수)다. 주어진 $(q,T,\{\xi_j\})$ 에서 이
식을 만족하는 $V_n$ 이 내부 전위다. 이는 다공성 전극 이론의 implicit 화학퍼텐셜 관계와 같은 구조다
\cite{doyle1993,newman}.

\textbf{배경 용량의 의미.} $Q_\bg(V_n,T)$ 는 단순 그래프 기저선이 아니라 뚜렷한 staging 으로 분리하지 않은
방전 방향 잔류 누적 용량이며, $\xi_j$ 가 평형 진행률을 즉시 따라가지 못할 때 전하 보존을 맞추기 위해 $V_n$
이 얼마나 이동해야 하는지를 정한다. 그 전위 기울기 $\partial Q_\bg/\partial V_n$ 이 잔류 chemical
capacitance 다. $V_n$ 을 안정적으로 풀려면 이 기울기가 양이어야 한다(\AL{5}; bazant2013 의 비평형
열역학적 미분용량):
\begin{equation}
\frac{\partial Q_\bg}{\partial V_n}\ \ge\ \epsilon_Q\ >\ 0 .
\label{eq:monotone}
\end{equation}
$\epsilon_Q$ 는 임의의 정규화 상수가 아니라 양의 미분용량이라는 물리 조건이며, 수치적으로는 해의 특이점을
막는 하한이다(\AL{5} BOUNDED).

\textbf{평형 극한 = OCV.} 평형(또는 충분히 느린 조건)에서는 $\xi_j=\xi_{\eq,j}(V_n,T)$ 이고
식~\eqref{eq:charge_balance} 가 $Q_\cell q=Q_\bg+\sum_j Q_{j,\tot}\xi_{\eq,j}$ 가 되어 그 해가 평형 음극
전위 $V_{n,\OCV}(q,T)$ 다. \textbf{즉 OCV 곡선은 임의로 주어진 lookup 이 아니라 전하 보존식의 평형
특수해다} — 중심식이 ``OCV 를 읽는'' 흐름으로 회귀하지 않는다.

\textbf{해 존재 조건(well-posedness).} 식~\eqref{eq:charge_balance} 가 $V_n$ 해를 가지려면 우변 목표가 배경
용량의 치역 안에 있어야 한다(\AL{7}): $Q_\cell q-\sum_j Q_{j,\tot}\xi_j\in\mathrm{Range}[Q_\bg(\cdot,T)]$.
단조성(식~\eqref{eq:monotone})과 이 치역 조건을 함께 만족하면 중간값 정리로 해가 유일하게 존재한다(부분
구간 피팅용 상대 기준형$\cdot$총용량 정합 등 실무는 \S\ref{sec:ch1_simplefit}).

% ====================================================================
\section{무대 (II): 평형 진행률 $\xi_\eq$ (logistic)}\label{sec:eq}
% ====================================================================
\textbf{왜 logistic 인가 (무비약 유도).} 흑연 intercalation 의 한 전이를 격자기체(lattice-gas)/정규용액
(regular-solution) 으로 보면, 점유율 $\xi$ 인 상태의 화학퍼텐셜은 \cite{mckinnon1983,bazant2013}(\AL{2})
\begin{equation}
\mu(\xi)=\mu^0+RT\ln\frac{\xi}{1-\xi}+\Omega_j(1-2\xi),
\label{eq:mu_latticegas}
\end{equation}
두 항의 출처를 학부 통계역학으로 한 줄씩 확인한다(식~\eqref{eq:mu_latticegas} 를 인용으로만 받지 않고 본문에서 유도한다).

\textbf{(i) 혼합 엔트로피 항.} $N$ 개의 동등한 자리 중 $n$ 개가 점유된 배열의 가짓수는
$W=\dfrac{N!}{n!\,(N-n)!}$ 이며, 점유율 $\xi\equiv n/N$ 으로 둔다. Boltzmann 관계 $S=k_B\ln W$ 에 Stirling
근사 $\ln m!\simeq m\ln m-m$ 를 적용하면(세 계승의 $-m$ 선형항이 $N=n+(N-n)$ 으로 상쇄되어)
\begin{equation}
S_{\mathrm{mix}}=k_B\big[N\ln N-n\ln n-(N-n)\ln(N-n)\big]
=-k_B N\big[\xi\ln\xi+(1-\xi)\ln(1-\xi)\big],
\label{eq:smix}
\end{equation}
즉 몰 기준($N\!\to\!N_A$, $R=k_B N_A$)으로 $S_{\mathrm{mix}}=-R[\xi\ln\xi+(1-\xi)\ln(1-\xi)]$ 다. 자유
에너지의 혼합 기여 $G_{\mathrm{mix}}=-TS_{\mathrm{mix}}=RT[\xi\ln\xi+(1-\xi)\ln(1-\xi)]$ 를 점유율로 미분하면
\begin{equation}
\frac{\partial G_{\mathrm{mix}}}{\partial\xi}=RT\big[(\ln\xi+1)-(\ln(1-\xi)+1)\big]=RT\ln\frac{\xi}{1-\xi}
\label{eq:mu_mix}
\end{equation}
($+1$ 과 $-1$ 이 상쇄). 여기서 $G_{\mathrm{mix}}$ 는 점유 자리 1몰당(intensive) 자유에너지이고, 화학퍼텐셜의 정의는
$\mu_{\mathrm{config}}=\partial G^{\mathrm{tot}}/\partial n$ ($G^{\mathrm{tot}}=N G_{\mathrm{mix}}$, 자리수 $N$, $\xi=n/N$)
인데 $\partial/\partial n=(1/N)\,\partial/\partial\xi$ 라 $N$ 인자가 상쇄되어 $\mu_{\mathrm{config}}=\partial G_{\mathrm{mix}}/\partial\xi$
다. 따라서 위 ξ-미분이 곧 화학퍼텐셜의 혼합 기여이며, 식~\eqref{eq:mu_latticegas} 의 첫 로그 항이 정확히 재현된다.

\textbf{(ii) 상호작용 항.} 정규용액 근사에서 점유 입자 간 상호작용(혼합 엔탈피)은
$G_{\mathrm{int}}=\Omega_j\,\xi(1-\xi)$ 이고, 그 미분 $\partial G_{\mathrm{int}}/\partial\xi
=\Omega_j(1-2\xi)$ (곱미분 $\dd[\xi-\xi^2]/\dd\xi=1-2\xi$)가 둘째 항이다. 표준 상태 항 $\mu^0$ 까지 더하면
$\mu=\mu^0+RT\ln[\xi/(1-\xi)]+\Omega_j(1-2\xi)$ 로 식~\eqref{eq:mu_latticegas} 가 완성된다(\AL{2};
$\mu^0,\Omega_j$ 의 미시값은 \cite{mckinnon1983,bazant2013} 에 위임).

\textbf{(iii) 전기화학 평형.} 평형은 이 화학퍼텐셜이 전극 전위와 균형을 이루는 조건
$\mu(\xi_\eq)=s_{\phi,j}F(V_n-U_j^0)$ 이다($s_{\phi,j}$=방전 방향 부호 인자, \S\ref{sec:notation};
eq:Aj 와 동일 convention). ideal 극한($\Omega_j=0$)에서 $\mu^0+RT\ln[\xi_\eq/(1-\xi_\eq)]
=s_{\phi,j}F(V_n-U_j^0)$ 이고, 상수 $\mu^0$ 를 기준 전위에 흡수해 $U_j\equiv U_j^0-\mu^0/(s_{\phi,j}F)$ 로
재정의하면
\begin{equation}
RT\ln\frac{\xi_\eq}{1-\xi_\eq}=s_{\phi,j}F(V_n-U_j).
\label{eq:eqbal_ideal}
\end{equation}

\textbf{(iv) $\xi_\eq$ 에 대해 풀기(한 줄씩).} 식~\eqref{eq:eqbal_ideal} 양변을 $RT$ 로 나누고 폭
$w_j\equiv RT/F$ 를 정의하면 $\ln[\xi_\eq/(1-\xi_\eq)]=s_{\phi,j}(V_n-U_j)/w_j$. 양변에 지수를 취해
$\dfrac{\xi_\eq}{1-\xi_\eq}=E$, $E\equiv\exp[s_{\phi,j}(V_n-U_j)/w_j]$. 이를 풀면 $\xi_\eq=E(1-\xi_\eq)$
$\Rightarrow\xi_\eq(1+E)=E\Rightarrow\xi_\eq=E/(1+E)$, 분자·분모를 $E$ 로 나눠 $1/E=\exp[-s_{\phi,j}(V_n-U_j)/w_j]$
이므로
\begin{equation}
\boxed{\;\xi_{\eq,j}(V_n,T)=\frac{1}{1+\exp[-s_{\phi,j}(V_n-U_j(T))/w_j(T)]}\;,\qquad w_j=\frac{RT}{F}\ (\text{ideal})\;.}
\label{eq:logistic}
\end{equation}
모든 중간 단계가 위에 명시됐다(방전 기준 $s_{\phi,j}=+1$ 이 기본이며, 이때 $V_n$ 증가가 $\xi_\eq$ 를 증가시킨다).
$\Omega_j\ne0$ 이면 식~\eqref{eq:eqbal_ideal} 우변에 $-\Omega_j(1-2\xi_\eq)$ 가 더해져 \emph{음함수}(implicit
isotherm)가 되며 단순 logistic 이 아니다 — 이때 $w_j$ 를 effective width 로 두는 것은 닫힌 유도가 아니라
coarse-grained ansatz 임을 명시한다(\AL{2} BOUNDED; $\Omega_j$ 가 크면 다중해·상분리 가능). 이는 매끄러운 logistic 곡선이며 $0\!\to\!1$ 급점프가
아니다. $\Omega_j\ne0$ 이면 폭·형태가 매끄럽게 보정되며, 실제로는 nonideality$\cdot$도메인 불균일$\cdot$유한
전이폭을 흡수하는 effective width $w_j$ 로 두어 $RT/F$ 에 강제로 묶지 않는다(\AL{4}). 이때
$n_j^{\eff}\equiv RT/(Fw_j)$ 가 effective 전자수 역할을 하며 Ch.3$\cdot$Ch.4 로 전달된다.

\textbf{저온 거동(서 단계 (1) 의 근거; 무비약).} logistic 의 peak 형상은 그 도함수로 결정된다. 식~
\eqref{eq:logistic} 에서 $z\equiv s_{\phi,j}(V_n-U_j)/w_j$ 로 두면 $\xi_\eq=(1+e^{-z})^{-1}$ 이고
$\dd z/\dd V_n=s_{\phi,j}/w_j$. 연쇄법칙으로 $\partial\xi_\eq/\partial V_n=(\dd\xi_\eq/\dd z)(\dd z/\dd V_n)$
이며, $\dd\xi_\eq/\dd z=-(1+e^{-z})^{-2}\cdot(-e^{-z})=e^{-z}/(1+e^{-z})^2$ 다. 여기서
$1-\xi_\eq=e^{-z}/(1+e^{-z})$ 이므로 $e^{-z}/(1+e^{-z})^2=[1/(1+e^{-z})]\cdot[e^{-z}/(1+e^{-z})]
=\xi_\eq(1-\xi_\eq)$ (표준 logistic 항등식)이다. 따라서
\begin{equation}
\frac{\partial\xi_{\eq,j}}{\partial V_n}=\frac{s_{\phi,j}\,\xi_{\eq,j}(1-\xi_{\eq,j})}{w_j(T)}
\label{eq:dxidV}
\end{equation}
로 결정되어 peak 높이 $\propto 1/w_j$, 폭 $\propto w_j$ 다($s_{\phi,j}=+1$ 방전 기준). 저온이면 \emph{ideal}
$w_j=RT/F$ 가 \emph{감소}하므로
peak 가 더 좁고 가팔라진다 — 즉 꼬리가 더 빨리 끝난다. 수치로 $25^\circ$C 에서 $w=RT/F\approx25.7$ mV,
$-15^\circ$C 에서 $\approx22.2$ mV($8.314\times298/96485=25.68$ mV, $\times258/96485=22.23$ mV).
\textbf{따라서 평형 열역학의 예측은 ``저온 $\to$ 짧은 꼬리''} 이고, 관찰(저온 긴 꼬리)은 평형 broadening 으로
설명할 수 없다. 꼬리의 주역은 평형이 아니라 동역학이어야 한다(\S\ref{sec:rate}--\ref{sec:dyn}).
\begin{flagbox}
\textbf{평형 형태는 데이터가 정한다.} logistic 대신 erf/Gaussian-cumulative isotherm 을 쓰는 것은 흑연에
대한 열역학적 유도가 문헌에 없고(무질서 반도체 Gaussian-disorder 모형 \cite{baessler1993} 과의 유추 +
경험적 peak-fit 수준) FLAGGED ansatz 다. 실제 형태는 저속 OCV 의 near-peak 감쇠를 $\log(dQ/dV)$ 대
$(V-U)$(지수=logistic) 또는 $(V-U)^2$(가우시안)로 보고 판별한다(\S\ref{sec:falsify}). \textbf{본 장의
꼬리 결론은 평형 형태에 무관}하다 — post-peak 에서 $\dd\xi_\eq/\dd q\simeq0$ 만 쓰기 때문이다.
\end{flagbox}

% ====================================================================
\section{주역 (I): 극판 전위에 의한 배리어 낮춤과 속도상수}\label{sec:rate}
% ====================================================================
이 절이 본 장의 핵심이다 — \emph{극판 전위가 어떻게 상변이 속도를 바꾸는가}.

\textbf{세 전위의 구분.} 측정되는 apparent 전위는 내부 전위에 분극이 더해진 값이다(\AL{6}):
\begin{equation}
V_{n,\app}=V_n+s_I|I|R_n(q,T,|I|).
\label{eq:Vapp}
\end{equation}
상변이 속도를 정하는 구동 전위는 $V_{n,\drive}$ 이며 가장 보수적인 기본형은 $V_{n,\drive}=V_n$ 이다(국소
계면 과전압을 분리하는 일반형 $V_{n,\drive}=V_n+\eta_{n,\drive}$ 는 Ch.3). \textbf{$R_n$(전압축 이동)과
아래의 속도 보조항을 같은 데이터로 동시에 자유 피팅하지 않는다} — 역할이 다르다($R_n$=전압 강하,
$\chi_j$=상변이 가속). 이 구분이 무너지면 둘이 서로를 흉내 내어 식별이 불가능해진다(\S\ref{sec:falsify}).

\textbf{구동력.} 한 전이의 구동력은 구동 전위가 전이 중심 전위에서 벗어난 정도다(\AL{11}):
\begin{equation}
\mathcal A_j=s_{\phi,j}F\,[V_{n,\drive}-U_j(T)].
\label{eq:Aj}
\end{equation}
$F[\mathrm V]$ 가 $\mathrm{J/mol}$ 차원이므로 $\mathcal A_j$ 는 몰당 자유에너지 구동력이다. 부호 인자
$s_{\phi,j}$ 는 선택한 방전 방향 전이의 구동력 부호를 정한다: $\chi_j\ge0$ 이고 $s_{\phi,j}[V_{n,\drive}-U_j]>0$
일 때 방전 방향 전이가 촉진된다. (본 장은 ideal 폭 $n_j^{\eff}=RT/(Fw_j)=1$ 을 기본으로 두므로 위 형태를 쓴다;
일반형 $\mathcal A_j=s_{\phi,j}\,n_j^{\eff}F(V_{n,\drive}-U_j)$ 는 effective width 와 함께 Ch.3 에서 쓴다.)

\textbf{유효 배리어(무비약).} 고유 활성화 자유에너지는 전이상태 이론(\AL{10}, \cite{eyring1935})으로
$\Delta G_{a,j}(T)=\Delta H_{a,j}-T\Delta S_{a,j}$ 다. 전위 구동력이 이 배리어를 \emph{선형으로} 낮춘다 —
이것이 Brønsted--Evans--Polanyi/Butler--Volmer 류의 표준 자유에너지 관계다(\AL{11},
\cite{evanspolanyi1938,bardfaulkner,bazant2013}).

\emph{선형의 출처(소구동력 1차 Taylor; 무생략).} Marcus 이론에서 구동력 $\mathcal A_j$ 가 주어졌을 때
활성화 장벽은 포물면 교차로부터 $\Delta G^\ddagger(\mathcal A_j)=\dfrac{(\lambda-\mathcal A_j)^2}{4\lambda}
=\dfrac{\lambda}{4}-\dfrac{\mathcal A_j}{2}+\dfrac{\mathcal A_j^2}{4\lambda}$ ($\lambda$=재구성 에너지)다.
$\mathcal A_j=0$ 둘레 1차 Taylor 전개를 하면 상수항 $\Delta G^\ddagger(0)=\lambda/4$ 와 1차 항
$\big[\dd\Delta G^\ddagger/\dd\mathcal A_j\big]_{0}\mathcal A_j=-\tfrac12\mathcal A_j$ 가 남고, 2차 항
$\mathcal A_j^2/(4\lambda)$ 은 $|\mathcal A_j|\ll\lambda$ 에서 무시된다. 따라서 1차 계수의 크기가 전달 계수
$\chi_j$ 다. \emph{대칭 Marcus 포물면(forward/backward 동일 곡률 $\lambda$)의 1차 전개는 정확히 $\chi_j=1/2$
단일값만 준다}; $\chi_j\ne1/2$(일반 $0\le\chi_j\le1$)은 1차 Taylor 자체에서 나오는 것이 아니라 \emph{곡률 비대칭}
($\lambda_f\ne\lambda_b$) 또는 BEP/Butler--Volmer 선형 자유에너지 관계의 경험적 기울기에서 오며
(\AL{11}, \cite{evanspolanyi1938,bardfaulkner}), 그 일반 범위로 둔다. 이 1차 항이 장벽을 낮추므로($\mathcal A_j>0$ 일 때 $-\chi_j\mathcal A_j<0$)
\begin{equation}
\boxed{\;\Delta G_{\eff,j}=\Delta G_{a,j}(T)-\chi_j\,\mathcal A_j\;=\;G-\chi_j\mathcal A_j\;}
\label{eq:Geff}
\end{equation}
($G\equiv\Delta G_{a,j}$ 를 mode 의 intrinsic barrier 로 표기, 위 $\lambda/4$ 등 상수항 흡수;
\S\ref{sec:spectrum} 에서 분포). 곡률(2차 항)이 중요해지는 $|\mathcal A_j|\sim\lambda$ 영역은 아래 Marcus
bound 로 분리한다. \emph{주의}: 이 $\chi_j\mathcal A_j$ 는 forward/backward 를 같은 배율로 키우는
\emph{방향 없는 common-mode mobility 가속}이며(keystone 참조), forward 장벽만 낮추는 표준 Butler--Volmer
와는 다르다 — 방향성 분해는 Ch.3 Level B 에서 다룬다. 구동력이
클수록($\mathcal A_j>0$) 유효 배리어가 낮아진다. 속도상수는 Eyring rate(\AL{10})
\begin{equation}
\boxed{\;k_j=k_0(T)\exp\!\Big[-\frac{\Delta G_{\eff,j}}{RT}\Big]\;,\qquad k_0(T)=\frac{k_BT}{h}\kappa(T)\;.}
\label{eq:kact}
\end{equation}
학부 수준에서는 Boltzmann 인자 $k_j\propto e^{-\Delta G_{\eff,j}/RT}$ 만으로 이후 논증에 충분하다 — 본 장의
모든 결론이 속도의 \emph{비}로 정의되는 길이 $L=|I|/(Q_\cell k_j)$(\S\ref{sec:dyn})로만 쓰여 prefactor 절대값
$k_0$ 가 상쇄되기 때문이다. prefactor 가 $k_BT/h$ 로 \emph{고정}(임의 상수 아님)이라는 절대값 자체는 전이상태
이론(\cite{eyring1935}) 결과로 \emph{인용}한다. $\Delta G_{\eff,j}\le0$
(강한 전위 구동)은 선형 근사가 강구동 영역에 들어간 것으로 표시하되 $\max(\cdot,0)$ 같은 절단으로 자르지
\emph{않는다}(GS-4; 절단이 필요해 보이면 그것은 아래 물리적 병목으로 다룬다).
\begin{keybox}
\textbf{$\chi_j\mathcal A_j$ 는 ``방향 있는 배리어 낮춤''이 아니라 ``방향 없는 완화 속도(mobility) 가속''이다
(Level A).} 무비약 근거: 2-상태 반응속도론에서 net rate $\dot\xi_j=r_j^+(1-\xi_j)-r_j^-\xi_j$. 정상점은
$\dot\xi_j=0$ 에서 $r_j^+(1-\xi_{\eq,j})=r_j^-\xi_{\eq,j}\Rightarrow r_j^+=(r_j^++r_j^-)\xi_{\eq,j}
\Rightarrow\xi_{\eq,j}=r_j^+/(r_j^++r_j^-)$ 다. $r_j^\pm$ 를 $\xi_j$ 의 작은 변화에 대해 국소 상수로 두면
net rate 를 직접 재배열할 수 있다: $\dot\xi_j=r_j^+-r_j^+\xi_j-r_j^-\xi_j=r_j^+-(r_j^++r_j^-)\xi_j
=(r_j^++r_j^-)\big[r_j^+/(r_j^++r_j^-)-\xi_j\big]=(r_j^++r_j^-)(\xi_{\eq,j}-\xi_j)$. 따라서
$\dot\xi_j=k_j(\xi_{\eq,j}-\xi_j)$, $k_j=r_j^++r_j^-$ 가 인수분해로 동시에 나온다(\AL{12}). \emph{강조}: 비약은
``$r_j^\pm$ 를 ($\xi_j$ 에 무관한) 국소 상수로 본다''는 \emph{가정}에만 있고, 그 가정 위에서는 위 재배열이 Taylor
전개가 아니라 대수적 항등이다(반대로 $r_j^\pm(\xi_j)$ 면 같은 식이 국소 선형화가 된다). 즉 $k_j$=완화 속도(목표에 접근하는 빠르기),
$\xi_{\eq,j}$=목표(방향). 식~\eqref{eq:kact} 처럼 $\chi_j\mathcal A_j$ 가 \emph{$k_j$
전체}에 들어가면 $r_j^+,r_j^-$ 가 같은 배율로 커져 비 $r_j^+/r_j^-=\xi_\eq/(1-\xi_\eq)$ 가 \emph{불변}
— 즉 방향(목표)은 열역학(\S\ref{sec:eq})이 전담하고, $\chi_j\mathcal A_j$ 는 방향 없는 속도 scale 일 뿐이다.
\emph{방향을 가진} 배리어 낮춤(forward 만 $-\beta\mathcal A$,
backward 는 $+(1-\beta)\mathcal A$ 라 비가 바뀜)은 \textbf{Ch.3 의 Level B} 에서 도입하며, 그때 본 장의
$\chi_j$ 가 transfer coefficient $\beta_j$ 와 같아진다(\AL{11}). \textbf{한정}: 본 장 Level A 의 ``공통 배율''
(forward/backward 동일 가속)은 Marcus 의 \emph{비대칭} 장벽 이동($\lambda\mp\mathcal A_j$)을 대칭으로 평균낸
\emph{모형 선택}이며, 이 대칭 극한에서만 forward/backward 가속이 세부 균형(detailed balance) 비를 보존한다. Level B(Ch.3)와
일치하는 것은 \emph{정상 target 극한} $\xi_{\mathrm{ss},j}\to\xi_{\eq,j}$ 한정이다(CHARTER step4). 본 장
결론식에서는 ``배리어 낮춤''이라는 방향성 함의 용어를 피한다.
\end{keybox}
\begin{boundbox}
\textbf{Marcus bound(\AL{15}).} 선형 $-\chi_j\mathcal A_j$ 는 작은 구동력의 1차 근사다. Marcus 이론
\cite{marcus1956} 은 곡률을 주므로 $|\mathcal A_j|$ 가 클 때(역전 영역 포함) 깨진다. 본 선형식은 꼬리 영역
$V_{n,\drive}\approx U_j$ 에서 유효하며 전역 정당화가 아니다.
\end{boundbox}

\textbf{물리적 병목(속도 절단이 아니라 직렬 합성).} 전하전달이 빨라져도 다음 단계(전해질/고체 수송, 유한
활성 자리, 고체상 확산, 유한 전류)가 전체 진행을 제한할 수 있다. 이를 \emph{시간척도의 직렬 합}으로 둔다
(수치 cap 이 아니라 율속 단계 합성):
\begin{equation}
\frac{1}{k_j^{\eff}}=\frac{1}{k_j}+\frac{1}{k_{j,\lim}},
\label{eq:kphys}
\end{equation}
$k_{j,\lim}\to\infty$ 면 활성화 지배($k_j^{\eff}\to k_j$), 반대면 병목 지배로 매끄럽게 환원된다.\footnote{%
$k_{j,\lim}$ 은 수송$\cdot$유한 자리$\cdot$고체 확산$\cdot$유한 전류 등 병목의 합($1/k_{j,\lim}=\sum 1/k_{j,m}$)
이다. 독립 자료 없이 자유 파라미터로 추가하면 $k_0,\Delta G_a,\chi_j$ 와 강하게 상관되므로 물리적 근거와
독립 제약이 있을 때만 쓴다. 유한 전류를 매끄럽지 않은 clamp($\min$/step)로 강제하지 않는다(GS-4); 그런
편의 처리와 피팅 안정화는 부록 \S\ref{sec:ch1_numeric} 에서 다룬다.} 이하 $k_j$ 는 이 $k_j^{\eff}$ 를 가리킨다.

% ====================================================================
\section{주역 (II): 진행률 동역학과 single-mode kernel}\label{sec:dyn}
% ====================================================================
상변이는 평형 진행률을 향해 유한 속도로 완화된다(\AL{12}, \cite{onsager1931,degrootmazur1962}):
\begin{equation}
\boxed{\;\frac{\dd\xi_j}{\dd t}=k_j\,[\xi_{\eq,j}(V_n,T)-\xi_j]\;.}
\label{eq:dxidt}
\end{equation}
$\xi_j$ 가 평형을 즉시 따라가지 못하면 전압 곡선과 ICA/DVA 에 꼬리가 생긴다. 정전류 구간($|I|>0$)에서는
$\dd q/\dd t=|I|/Q_\cell$ 이므로 그 역수가 $\dd t/\dd q=Q_\cell/|I|$ 이고, 연쇄법칙
$\dd\xi_j/\dd q=(\dd\xi_j/\dd t)(\dd t/\dd q)$ 에 식~\eqref{eq:dxidt} 의 $\dd\xi_j/\dd t=k_j[\xi_{\eq,j}-\xi_j]$
를 넣으면
\begin{equation}
\frac{\dd\xi_j}{\dd q}=\frac{Q_\cell}{|I|}\,k_j\,[\xi_{\eq,j}-\xi_j].
\label{eq:dxidq}
\end{equation}

\textbf{Single-mode kernel.} 한 mode 의 지연을 $r=\xi_{\eq,j}-\xi_j$ 라 하자. peak 를 지난 영역(post-peak)
에서 목표가 거의 정지하고($\dd\xi_{\eq,j}/\dd q\simeq0$) 구동력도 완만하면($\dd\mathcal A_j/\dd q\simeq0$,
즉 $V_n$ 완만 — \emph{별개 가정 두 개}임에 유의), 식~\eqref{eq:dxidq} 는 1계 선형 ODE 가 된다.
\emph{무비약 유도}: 정의 $r=\xi_{\eq,j}-\xi_j$ 를 $q$ 로 미분하면 $\dd r/\dd q=\dd\xi_{\eq,j}/\dd q-\dd\xi_j/\dd q$
이고, post-peak 가정 $\dd\xi_{\eq,j}/\dd q\simeq0$ 으로 첫 항이 사라져 $\dd r/\dd q=-\dd\xi_j/\dd q$ 다. 여기에
식~\eqref{eq:dxidq} $\dd\xi_j/\dd q=(Q_\cell k_j/|I|)(\xi_{\eq,j}-\xi_j)=(Q_\cell k_j/|I|)\,r$ 를 넣고
$L\equiv|I|/(Q_\cell k_j)$ 로 두면 $\dd r/\dd q=-(1/L)\,r$ 다. 여기서 $L$ 은 일반적으로 $q$ 의 함수다($k_j$ 가
구동력 $\mathcal A_j$ 에 지수 의존, 식~\eqref{eq:kact}). \emph{그러나} 위 가정 (나) $\dd\mathcal A_j/\dd q\simeq0$
이면 $k_j$, 따라서 $L$ 이 이 구간에서 $q$-무관(상수)이므로 $L$ 을 적분 밖으로 뺄 수 있다 — 이것이 ``$L$ 일정''의
근거다(별개 가정이 아니라 (나)의 \emph{귀결}). 변수분리 $\dd r/r=-\dd q/L$ 를 $q_a$ 에서 $q$ 까지 적분하면
$\int_{q_a}^q \dd q'/L=(q-q_a)/L$, 즉 $\ln[r(q)/r(q_a)]=-(q-q_a)/L$ 이고 양변에 지수를 취해 지수 감쇠가 따라온다:
\begin{equation}
r(q)\simeq r(q_a)\exp\!\Big[-\frac{q-q_a}{L}\Big],\qquad \boxed{\;L=\frac{|I|}{Q_\cell\,k_j}\;}
\label{eq:single_kernel}
\end{equation}
(차원: $\mathrm A=\mathrm{C/s}$, 속도상수 $k$ 는 $\mathrm{s^{-1}}$, $Q_\cell$ 은 $\mathrm C$ 이므로
$|I|/(Q_\cell k)=(\mathrm{C/s})/(\mathrm C\cdot\mathrm{s^{-1}})=(\mathrm{C/s})/(\mathrm{C/s})=1$, 즉 무차원).
이것은 \emph{관측 꼬리 전체가 단일
지수}라는 뜻이 아니라 \textbf{하나의 mode 가 만드는 지수 kernel} 이다. 관측 꼬리는 여러 mode 의 중첩이다
(\S\ref{sec:kernel}).

\textbf{잔차에서 진행률 도함수로(무비약 연결).} 다음 절의 중첩은 잔차 $r$ 이 아니라 \emph{진행률 도함수}
$\dd\xi_j/\dd q$ 의 합이다. 둘의 연결은 식~\eqref{eq:dxidq} 에서 $\dd\xi_j/\dd q=(Q_\cell k_j/|I|)\,r=r/L$
이므로
\begin{equation}
\frac{\dd\xi_j}{\dd q}=\frac{r(q_a)}{L}\exp\!\Big[-\frac{q-q_a}{L}\Big].
\label{eq:single_dxidq}
\end{equation}
바로 이 \textbf{$1/L$ 인자}가 \S\ref{sec:kernel} kernel 의 $1/L$ 출처다 — 중첩은 각 mode 의 \emph{진행률
기여}를 합한 것이지 잔차를 합한 것이 아니다.

\textbf{저온 거동(서 단계 (2)).} $k_j=k_0e^{-(G-\chi\mathcal A)/RT}$ 가 저온서 작아져 $L$ 이 커진다 $\Rightarrow$
꼬리가 길어진다 — 관찰과 일치. 빠른 kinetics 극한($k_j\to\infty$)에서는 $L\to0$ 이며 $\dd\xi_j/\dd q\to
\dd\xi_{\eq,j}/\dd q$ 로 접근한다(작은 지연 $\times$ 큰 속도의 곱이 유한; bracket 이 단순히 0 이 되는 것이
아니라 점근 균형의 결과다).

% ====================================================================
\section{주역 (III): 배리어 분포 $\to$ relaxation-length spectrum}\label{sec:spectrum}
% ====================================================================
\textbf{물리.} 실제 전극은 단일 배리어 $G$ 가 아니라 입자 크기$\cdot$국소 staging 환경$\cdot$결정립계$\cdot$결함
$\cdot$응력 차이로 인한 \emph{배리어 분포} $\rho_G(G;T)$ 를 가진다(\AL{13}). 진행률을 특정 배리어 기준으로
자르지 않고, 배리어 $g$ 를 가진 집단의 진행률을 $\xi_j(g,t)$ 로 두어 분포로 평균한다($\int_0^\infty
\rho_G(g;T)\,\dd g=1$):
\begin{equation}
\xi_j(t)=\int_0^\infty \rho_G(g;T)\,\xi_j(g,t)\,\dd g,\qquad
\frac{\dd\xi_j(g,t)}{\dd t}=k_j(g)\,[\xi_{\eq,j}-\xi_j(g,t)].
\label{eq:xi_dist}
\end{equation}
\begin{flagbox}
\textbf{독립 완화 가정(평균장; \AL{13}).} 식~\eqref{eq:xi_dist} 는 배리어 $g$ 별 집단이 \emph{서로 독립적으로}
1차 완화한다고 본다(mode 간 결합 무시). 그러나 흑연 staging 은 nucleation$\cdot$상경계 이동을 동반하는
\emph{협동적 1차 상전이}다(\AL{3}, \cite{funabiki1999ea,funabiki1999jes}). 따라서 ``독립 평행 1차 완화의 분포''
근사는 mode 간 상호작용(상경계 결합$\cdot$Avrami 형 비선형 동역학)을 무시한 \emph{평균장 ansatz} 이며, 그 타당성은
FLAGGED — Ch.3(반응속도론 일반화)$\cdot$본 장 부록 \S\ref{sec:ch1_numeric}(수치)에서 검증한다.
\end{flagbox}

\textbf{배리어 $\to$ 길이 지수 매핑(꼬리가 늘어지는 근원).} 관측 꼬리의 모양은 $\rho_G$ 와 \emph{같지 않다}
— 배리어 $G$ 가 먼저 속도로, 다시 길이로 \emph{지수 변환}되기 때문이다. 식~\eqref{eq:single_kernel} 의
$L=|I|/(Q_\cell k_j)$ 에 식~\eqref{eq:kact} $k_j=k_0(T)\exp[-(G-\chi_j\mathcal A_j)/RT]$(mode 별 intrinsic
barrier $G\equiv\Delta G_{a,j}$, $\Delta G_{\eff,j}=G-\chi_j\mathcal A_j$)를 대입한다. \emph{이 닫힌 지수 매핑은
활성화 지배 극한}($k_j^{\eff}\simeq k_j$, 즉 식~\eqref{eq:kphys} 의 $k_{j,\lim}\to\infty$)\emph{을 전제}한다 — 꼬리
영역($V_{n,\drive}\approx U_j$, 저구동)에서는 활성화가 율속이라 정당하나, 병목이 유효하면
$1/k_j^{\eff}=1/k_j+1/k_{j,\lim}$ 가 지수 밖 유리식으로 들어가 $L(G)$ 가 단순 지수형이 아니다(그 경우의 수치 풀이는
부록 \S\ref{sec:ch1_numeric}). $L$ 은 $k_j$ 의 \emph{역수}에 비례하므로 $1/k_j=(1/k_0)\exp[+(G-\chi_j\mathcal A_j)/RT]$ 로
\emph{지수 부호가 반전}되어(속도가 작을수록 길이가 길다)
\begin{equation}
\boxed{\;L(G,T,V_{n,\drive})=\frac{|I|}{Q_\cell\,k_0(T)}\exp\!\Big[\frac{G-\chi_j\mathcal A_j}{RT}\Big]\;.}
\label{eq:L_of_G}
\end{equation}
양변에 로그를 취하면 $\ln L=\ln L_0+(G-\chi_j\mathcal A_j)/RT$ 이다. 즉 배리어의 표준편차 $\sigma_G$($\rho_G$ 의
표준편차)가 $\ln L$ 에서 $\sigma_G/RT$ 로 나타나, \textbf{저온일수록 $\ln L$ 의 퍼짐이 커진다} — 작은 배리어 분산이 길이에서
지수적으로 확대된다. 이것이 ``저온 긴 꼬리''의 근원이며, 고정된 배리어 분포가 온도를 통해 stretched 거동을
만든다는 메커니즘은 무질서계$\cdot$빠른 이온 전도체 동역학에서 확립돼 있다(\AL{14},
\cite{svare2000,johnston2006,lindsey1980}; 단 graphite ICA 꼬리로의 \emph{적용}은 BOUNDED, \S\ref{sec:kernel}).

\textbf{변수변환: 배리어 분포 $\to$ 완화 길이 분포(무비약, 단계별).} 우리가 만들 변수 $A_L(L)$ 은 ``완화 길이가
$L$ 인 mode 가 꼬리에 \emph{얼마나 기여}하는가''의 분포다. 이는 \emph{이미 아는} 배리어 분포 $\rho_G(G)$ 를
식~\eqref{eq:L_of_G} 의 $L(G)$ 로 \emph{좌표만 바꾼} 것 — 새 물리를 넣는 게 아니라 같은 mode 집단을 길이 축에서
다시 세는 것이다. 세 단계로 만든다.

\emph{단계 1 — 자코비안(밀도 보존).} 확률(mode 분율)은 좌표를 바꿔도 보존되므로 $A_L^{\mathrm{prob}}(L)\,\dd L
=\rho_G(G)\,\dd G$. 역변환(식~\eqref{eq:L_of_G} 을 $G$ 에 대해 풀면) $G(L)=\chi_j\mathcal A_j+RT\ln(L/L_0)$
($L_0\equiv|I|/(Q_\cell k_0)$; $G=0$ 일 때 $L=L_{\min}$, 단계 2 에서 정식화)에서 $\dd G/\dd L=RT/L>0$ 이라
일대일이므로
\begin{equation}
A_L^{\mathrm{prob}}(L)=\rho_G\big(G(L)\big)\,\Big|\frac{\dd G}{\dd L}\Big|=\rho_G\big(G(L)\big)\,\frac{RT}{L}.
\label{eq:spectrum_prob}
\end{equation}
\emph{의미}: $A_L^{\mathrm{prob}}(L)\dd L$ $=$ 완화 길이가 $[L,L+\dd L]$ 인 mode 의 \emph{분율}, 규격 보존
($\int A_L^{\mathrm{prob}}\dd L=\int\rho_G\dd G=1$). 자코비안 $RT/L$ 은 ``배리어 한 칸 $=$ 길이 몇 칸''의 환산율일
뿐 임의 가중이 아니다.

\emph{단계 2 — 지지 범위(음의 배리어 금지).} 활성화 배리어는 음수가 못 되므로($G\ge0$, \S\ref{sec:notation})
식~\eqref{eq:L_of_G} 의 단조성상 $G\ge0\Leftrightarrow L\ge L_{\min}$, $L_{\min}\equiv L_0 e^{-\chi_j\mathcal A_j/RT}$
($G=0$ 대응 최단 길이)다. 따라서 $L<L_{\min}$ 에서 $A_L^{\mathrm{prob}}=0$ — 식에 Heaviside $H(L-L_{\min})$ 로
달아 음의 배리어 영역을 잘못 적분하지 않게 한다.

\emph{단계 3 — 개수 밀도 $\to$ 용량 가중(측도 변환, 무비약).} 단계 1$\cdot$2 의 $A_L^{\mathrm{prob}}$ 는 ``각 길이에
mode 가 몇이나''(개수 밀도)다. 그러나 관측 꼬리는 \emph{용량 가중} 진행 $\Theta=Q_p^{-1}\sum_m Q_m\xi_m$
(\S\ref{sec:fiteq}; $Q_m$=mode 의 용량 기여)로 나타나므로, 길이별 기여는 \emph{개수}$\times$\emph{mode 당 용량}이다.
이산 합을 연속 적분으로 옮기면
\begin{equation*}
\Theta=Q_p^{-1}\sum_m Q_m\,\xi_m\;\xrightarrow{\text{연속극한}}\;Q_p^{-1}\!\int Q(L)\,\xi(L)\,A_L^{\mathrm{prob}}(L)\,\dd L,
\end{equation*}
즉 개수 밀도 $A_L^{\mathrm{prob}}$ 에 \emph{mode 당 용량} $Q(L)$ 이 곱해진다. 이 용량 가중을 (무차원화한)
$A_0(L)\equiv Q(L)/\bar Q$ 로 두면 $A_0$ 는 개수 밀도 $\rho_G$ 에서 \emph{용량 측도}로 가는 Radon--Nikodym
가중이다 — 입자 크기 분포$\cdot$활성 site 밀도 같은 독립 물리량에서 오며 임의 피팅 knob 이 아니다. 관측 spectrum 은
그 곱이다:
\begin{equation}
\boxed{\;A_L(L)=A_0(L)\;\underbrace{\rho_G\big(G(L)\big)\,\frac{RT}{L}\,H(L-L_{\min})}_{\displaystyle A_L^{\mathrm{prob}}(L)\ (\text{규격 보존})}\;.}
\label{eq:spectrum}
\end{equation}
$A_0$ 는 \emph{임의 피팅 knob 이 아니라} 독립 물리량(입자 크기 분포 등)에서 와야 한다. $A_0\equiv1$ 이면
$A_L=A_L^{\mathrm{prob}}$ 로 규격 보존, $A_0\ne1$ 이면 $A_L$ 은 규격화 안 된 \emph{진폭 스펙트럼}이다 — \textbf{이하
kernel 은 이 $A_L$(진폭 spectrum)을 쓴다}(관측 꼬리는 진폭 가중 중첩이므로). 자코비안 $RT/L$ 과 아래 kernel 의
$1/L$ 인자는 \emph{출처가 다른 별개}임에 유의(\S\ref{sec:kernel}).
\begin{boundbox}
\textbf{비유일성(\AL{16}).} stretched 거동에서 $\rho_G$ 로의 역매핑은 \emph{비유일}하다(재포획 등 다른
메커니즘도 같은 꼬리를 낼 수 있다 \cite{baessler1993}). 따라서 $\rho_G$ 를 관측 꼬리에서 \emph{역산하지
않고}, 독립 측정/모형한 $\rho_G$ 로부터 꼬리를 \emph{forward 예측}만 한다(\S\ref{sec:falsify}).
\end{boundbox}

% ====================================================================
\section{관측 꼬리 = kernel integral}\label{sec:kernel}
% ====================================================================
\textbf{$A_L$ 의 물리 의미.} $A_L(L)\,\dd L$ 은 \emph{완화 길이가 $[L,L+\dd L]$ 인 mode 들이 전체 진행에 기여하는
분율}이다 — 즉 관측 꼬리를 ``완화 길이별로 분해한 분포''다. 빠른 mode(작은 $L$)는 peak 근처에서 빨리 끝나고,
느린 mode(큰 $L$)가 꼬리를 멀리 끌고 간다. 따라서 관측 꼬리는 \emph{서로 다른 $L$ 의 지수 감쇠를 $A_L$ 로 가중
합한 것}이다: 전체 진행률의 post-peak 꼬리는 각 mode 의 진행률 기여 $\dd\xi_j/\dd q=(r(q_a;L)/L)e^{-(q-q_a)/L}$
(식~\eqref{eq:single_dxidq})를 $A_L$ 로 가중 합한 것이다. 여기서 mode 별 초기 잔차 $r(q_a;L)$ 를 post-peak
공통값 $\bar r$ 로 \emph{근사}하면 $\bar r$ 가 적분 밖으로 빠지고, 이를 진폭 규격에 재흡수($A_L\!\leftarrow\!\bar r\,A_L$)해
아래 식을 얻는다 — $r(q_a)$ 가 사라진 게 아니라 $A_L$ 의 크기 인자로 흡수된 것이다:
\begin{equation}
\boxed{\;\frac{\dd\Theta_\mathrm{tail}}{\dd q}\simeq\int_0^\infty A_L(L)\,\frac1L\,
\exp\!\Big[-\frac{q-q_a}{L}\Big]\,\dd L\;.}
\label{eq:kernel_integral}
\end{equation}
이 균일 근사는 BOUNDED 다: 빠른 mode(작은 $L$)는 $q_a$ 에서 이미 더 진행해 $r(q_a;L)$ 가 작을 수 있어 원칙적으로
$L$-의존이다. $r(q_a;L)$ 이 $L$ 의 \emph{감소}함수이면 작은-$L$ 기여를 과대평가해 \emph{꼬리 앞부분이 실제보다
가팔라지는} 쪽으로 편향되며, 이 편향은 $r(q_a;L)$ 를 직접 보유하는 Plan B(g-grid, \S\ref{sec:planB})와 대조해
정량화한다. 피적분의 $1/L$ 인자는 진행률 도함수(식~\eqref{eq:single_dxidq})에서 온 것이고, spectrum 변환의
자코비안 $RT/L$ 은 $A_L$ 안에 이미 들어 있어 둘은 \emph{별개} 출처다 — 임의 가중이 아니다. $A_L$ 은 \S\ref{sec:spectrum}
의 $H(L-L_{\min})$ 를 품으므로 적분의 실효 하한은 $L_{\min}>0$ 이고 $1/L$ 의 $L\!\to\!0$ 발산은 일어나지 않는다.
이 적분은 relaxation-rate($1/L$) 분포 위의 Laplace 형 변환이며, 단일 지수의 \emph{연속 중첩}이 일반적으로 비지수
(stretched)$\sim$power-law 꼬리를 준다(\AL{14}, \cite{johnston2006}; KWW stretched 함수와 완화시간 분포의
등가성 \cite{lindsey1980}). single-mode(식~\eqref{eq:single_kernel})는 $A_L=\delta(L-L^\ast)$ 인 특수해다
($L^\ast$=그 mode 의 실제 완화 길이 $|I|/(Q_\cell k_j)$; \S\ref{sec:spectrum} 의 prefactor 길이 $L_0$ 와 구별).
또한 식~\eqref{eq:kernel_integral} 은 목표가 멈춘($\dd\xi_\eq/\dd q\!\simeq\!0$) post-peak \emph{자유감쇠} 극한의 꼬리이며,
목표가 움직이는 일반(자기참조 포함) 경우는 동일 kernel 을 \emph{목표} $\Theta_\eq$ 에 작용시킨 \S\ref{sec:volterra}
(식~\eqref{eq:volterra})로 일반화된다(잔차가 아니라 목표에 곱함 — \S\ref{sec:volterra} 의 이중계상 경고 참조).
\begin{boundbox}
\textbf{유효범위(\AL{14}, 정직 표기).} ``고정 배리어 분포 $\to$ 지수 매핑 $\to$ 저온 긴 꼬리''의 메커니즘은
빠른 이온 전도체$\cdot$무질서 완화에서 확립돼 있다(\cite{svare2000,johnston2006,lindsey1980}). 다만 (i)
특정 $\rho_G$ 형태가 \emph{정확히 어떤} stretched 지수를 주는지는 데이터/수치(부록 \S\ref{sec:ch1_numeric})로 정하며, (ii) 위
문헌은 전도도/유전 완화 맥락이라 graphite ICA 꼬리로의 \emph{적용}이 본 작업의 기여다. 이 두 항목은
established 가 아니라 BOUNDED 로 둔다.
\end{boundbox}

% ====================================================================
\section{자기참조 결합과 Volterra 적분식}\label{sec:volterra}
% ====================================================================
\textbf{왜 자기참조인가(물리).} $\xi_j$ 와 $V_n$ 은 전하 보존(식~\eqref{eq:charge_balance})으로 얽혀 있다 — 진행
$\Theta$ 가 늘면 $V_n$ 이 이동하고(식~\eqref{eq:charge_balance} 미분), $V_n$ 이 이동하면 평형 목표
$\Theta_\eq(V_n)$ 가 바뀌며, 바뀐 목표가 다시 $\Theta$ 의 완화 방향을 정한다. 이 되먹임 때문에 $\Theta$ 는 \emph{자기
자신을 목표로 갖는} 적분식의 해가 된다.

\textbf{적분형 유도(single-mode 에서, 무비약).} 한 mode 의 $q$-좌표 완화는 $\dd\Theta/\dd q=(1/L)(\Theta_\eq-\Theta)$
다(식~\eqref{eq:single_kernel} 와 같은 1차 완화, rate $1/L$). 이 1계 선형 ODE 를 상수변화법(적분인자 $e^{q/L}$)으로
풀면 $\dd(\Theta e^{q/L})/\dd q=(1/L)\Theta_\eq e^{q/L}$ 이고 $0$ 에서 $q$ 까지 적분해
\begin{equation}
\Theta(q)=\underbrace{\Theta(0)\,e^{-q/L}}_{\text{초기값 자유 감쇠}}+\int_0^q \frac1L\,e^{-(q-q')/L}\,\Theta_\eq(q')\,\dd q'.
\label{eq:singlemode_integral}
\end{equation}
즉 \textbf{kernel $(1/L)e^{-(q-q')/L}$ 이 \emph{목표} $\Theta_\eq$ 에 곱해진} 형이다(상수 목표 $\Theta_\eq=\Theta_0$,
$\Theta(0)=0$ 이면 $\Theta=\Theta_0(1-e^{-q/L})$ 로 복귀 — 무결성). 이제 두 곳을 일반화한다: \textbf{(i) 자기참조} —
목표가 미지 진행에 의존, $\Theta_\eq\to\Theta_\eq(V_n(q',\Theta(q')))$; \textbf{(ii) spectrum} — 단일 $L$ 을
분포 $A_L$ 로 적분. 목표 $\Theta_\eq$ 는 \emph{모든 mode 공통}이라(식~\eqref{eq:xi_dist} 의 $\xi_{\eq,j}$ 가
$g$-무관) $L$-적분 밖으로 빠지고 적분은 kernel 에만 작용하므로, $(1/L)e^{-(q-q')/L}\to\mathcal K(q,q')=\int
A_L(1/L)e^{-(q-q')/L}\dd L$, 초기항 $\Theta(0)e^{-q/L}\to\Theta_\mathrm{init}=\int A_L\Theta_{\mathrm{init},0}(L)e^{-q/L}\dd L$
($\Theta_{\mathrm{init},0}(L)$=mode 별 초기 진행, 상수 목표 scalar $\Theta_0$ 과 구별). 그러면 인과 상한 $q$ 의 Volterra 2종이
\emph{유도}된다(``뜬금없이''가 아니라 식~\eqref{eq:singlemode_integral} 의 두 일반화):
\begin{equation}
\Theta(q)=\Theta_\mathrm{init}(q)+\int_0^q \mathcal K(q,q')\,\Theta_\eq\big(V_n(q',\Theta(q')),T\big)\,\dd q',
\qquad \mathcal K(q,q')=\int_0^\infty A_L(L)\,\tfrac1L\,e^{-(q-q')/L}\,\dd L.
\label{eq:volterra}
\end{equation}
여기서 $\Theta_\mathrm{init}(q)=\int_0^\infty A_L(L)\,\Theta_{\mathrm{init},0}(L)\,e^{-q/L}\,\dd L$ 는 peak 진입 시점의 초기 mode
진행 $\Theta_{\mathrm{init},0}(L)$ 이 외부 구동 없이 자유 감쇠하는 동차항이고, 미지 $\Theta$ 는 좌변과 목표
$\Theta_\eq(V_n(q',\Theta(q')))$ 를 통해서만 등장한다 — 즉 ``현재 진행이 현재 목표를 바꾸고, 그 목표가 다시
현재 진행을 정한다''는 되먹임이 자기참조의 정체다.
\begin{boundbox}
\textbf{이중 계상 경고.} kernel $\mathcal K$ 는 반드시 \emph{목표} $\Theta_\eq$ 에 곱한다. 만약 잔차
$[\Theta_\eq-\Theta]$ 에 곱하면 상수 목표 극한에서 정상상태가 $\Theta_0/2$ 로 어긋난다(완화를 두 번 세는
오류). 상수 목표 $\Theta_\eq=\Theta_0$, $\mathcal K=(1/L)e^{-(q-q')/L}$, $\Theta_\mathrm{init}=0$ 극한에서
식~\eqref{eq:volterra} 는 $\Theta(q)=\Theta_0(1-e^{-q/L})$ 로 single-mode ODE(식~\eqref{eq:single_kernel})
와 정확히 일치하며 $q\to\infty$ 서 목표 $\Theta_0$ 로 수렴한다(single-mode 무결성 점검). \emph{분포} $A_L$ 의 경우
같은 극한의 $q\to\infty$ 값은 $\big(\int_0^\infty A_L\dd L\big)\Theta_0$ 이므로, $A_0\equiv1$(규격 보존, \S\ref{sec:spectrum})
일 때만 $\Theta_0$ 로 수렴한다 — $A_0\ne1$(진폭 spectrum)이면 총 scale 은 $\Theta\cdot Q_p$ 정의(\S\ref{sec:fiteq})의
용량 규격에 흡수되어 절대 scale 이 재조정된다(위 single-mode 점검은 $\delta$-collapse 한정).
\end{boundbox}
% ====================================================================
\section{자기참조식을 닫기 (1): Plan B — 직접 g-grid (항상 보존)}\label{sec:planB}
% ====================================================================
식~\eqref{eq:volterra} 에서 \emph{값}을 얻으려면 미지 $\Theta$ 를 닫아야 한다. 두 방법을 둔다(물리 결론은 어느
쪽으로 풀든 같다): \textbf{Plan B}(본 절, 항상 보존되는 수치 core/validator)와 \textbf{Plan A}(\S\ref{sec:planA},
검증되면 쓰는 닫힌형).

\textbf{Plan B — 직접 $g$-grid 수치 적분.} 배리어를 격자 $\{g_m\}$ 로 이산화해
식~\eqref{eq:xi_dist} 를 시간 적분하며, 매 시점 전하 보존(식~\eqref{eq:charge_balance})으로 $V_n$ 을 풀어 결합한다:
\begin{equation*}
\xi_j(t)=\sum_m w_m\,\rho_G(g_m)\,\xi_j(g_m,t),\qquad \sum_m w_m\,\rho_G(g_m)=\int_0^\infty\rho_G\,\dd g=1.
\end{equation*}
즉 $w_m$ 은 격자 $\{g_m\}$ 의 quadrature 가중으로 식~\eqref{eq:xi_dist} 의 \emph{적분을 그대로 이산화}한 것이지
새 가중을 넣은 게 아니며($w_m\rho_G(g_m)$ 이 정규화 측도를 이뤄 합이 1), 비균일$\cdot$large-$L$ 편중 격자
(부록 \S\ref{sec:ch1_numeric})에서도 이 정규화는 재계산해 보존한다. 주어진 $\rho_G$ 와 완화 closure 하에서
이산화 오차만 남는 수치 기준해이며, Plan A 닫힌형을 이 기준해와 상대오차
$\varepsilon=\lVert\Theta_\mathrm A-\Theta_\mathrm B\rVert/\lVert\Theta_\mathrm B\rVert$ 로 대조하는 \emph{validator}
를 겸한다(식~\eqref{eq:closed}, 수치 스킴$\cdot$수렴 판정 상세는 부록 \S\ref{sec:ch1_numeric}). 저온$\to$큰
$L$$\to$긴 꼬리 같은 \emph{물리 결론은 (위 완화 closure 전제 하에) 닫힌형 없이도 성립}한다.

% ====================================================================
\section{자기참조식을 닫기 (2): Plan A — 선형화 $\to$ 닫힌형}\label{sec:planA}
% ====================================================================
Plan B 가 항상 보존되는 수치해라면, Plan A 는 \emph{데이터 피팅에 직접 쓸} 닫힌 식이다(검증 통과 시). \textbf{baseline
$\Theta^{(0)}$ 정의(먼저 명시)}: 자기참조를 끈($V_n$ 고정, 아래 $b=0$) 0차 해를 $\Theta^{(0)}(q)$ 로 둔다 —
single-mode 극한에선 $\Theta^{(0)}(q)=\Theta_0(1-e^{-(q-q_a)/L})$(식~\eqref{eq:singlemode_integral} 의 상수목표
해), 일반적으로는 Volterra(식~\eqref{eq:volterra})를 $V_n$-고정으로 평가한 해다. Plan A 는 이 $\Theta^{(0)}$ 둘레의
self-consistency 보정이다.

\emph{단계 0(선형화).} 식~\eqref{eq:volterra} 의 비선형성은 오직 $\Theta_\eq(V_n(q',\Theta(q')))$ 한 곳이다.
$\Theta^{(0)}$ 둘레에서 1차 전개하면 $\Theta_\eq(V_n(\Theta))\approx a(q')+b(q')\,\Theta(q')$, 표준 1차 Taylor
분해상 절편은 $a(q')=\Theta_\eq\big(V_n(\Theta^{(0)}(q'))\big)-b(q')\,\Theta^{(0)}(q')$(baseline 에서 평가한 목표값에서
1차 기울기$\times$baseline 을 뺀 상수항)이고, 기울기(자기참조 계수)는
\begin{equation}
b(q')=\frac{\partial\Theta_\eq}{\partial V_n}\,\frac{\partial V_n}{\partial\Theta}
=-\frac{Q_p}{C_\bg}\,\frac{\partial\Theta_\eq}{\partial V_n}\ \le\ 0,
\label{eq:selfcoupling}
\end{equation}
$\partial V_n/\partial\Theta=-Q_p/C_\bg$ 는 전하 보존(식~\eqref{eq:charge_balance})을 $Q_p\Theta\equiv\sum_j
Q_{j,\tot}\xi_j$($Q_p=\sum_j Q_{j,\tot}$=총 phase 용량), $C_\bg\equiv\partial Q_\bg/\partial V_n$(배경 chemical
capacitance, 상세 \S\ref{sec:fiteq})로 묶은 $Q_\cell q=Q_\bg(V_n)+Q_p\Theta$ 를 $\Theta$ 로 미분해 얻고($C_\bg>0$), $\partial\Theta_\eq/\partial V_n>0$(평형 진행률은 전위에 증가, 식~\eqref{eq:dxidV} 의 $\dd\xi_\eq/\dd V_n>0$
를 $\Theta_\eq=Q_p^{-1}\sum_j Q_{j,\tot}\xi_{\eq,j}$ 로 합산)이므로 $b\le0$ 이 \emph{확정}된다. $b\le0$ 의 \emph{의미}는
``진행이 늘면 목표가 줄어 진행을 \emph{되끌어당긴다}''는 음의 되먹임이다.
이로써 식~\eqref{eq:volterra} 가 미지 $\Theta$ 에 \emph{선형}이 된다: $\Theta=c+\int_0^q\mathcal K\,b\,\Theta\,\dd q'$,
$c(q)=\Theta_\mathrm{init}+\int_0^q\mathcal K\,a\,\dd q'$.
\emph{단계 1(ratio ansatz).} 사용자 PhD \cite{lee2011,son2013} 의 ratio-substitution 을 \emph{미지} $\Theta(q')$
에만 적용한다: $\Theta(q')\approx\Theta(q)\,\Theta^{(0)}(q')/\Theta^{(0)}(q)$ (진행의 \emph{모양}은 baseline 을
따르고 크기만 미지로 둔다). \emph{정당화}: 자기참조항 $\int\mathcal K b\,\Theta$ 가 $c$ 에 대해 섭동($|b|$ 작음)일
때, 형상은 0차 baseline $\Theta^{(0)}$ 이 지배하고 $b$ 는 진폭만 재조정한다 — $b\to0$ 에서 정확(식~\eqref{eq:closed}
이 $c$ 로 복귀), $|b|$ 가 커질수록 형상 왜곡이 누적되어 아래 boundbox 의 부정확 한계로 이어진다.
\emph{단계 2($\Theta(q)$ 인수분해).} 단계 1 을 단계 0 의 선형식 $\Theta=c+\int_0^q\mathcal K b\,\Theta\,\dd q'$ 의
적분에 넣으면, 앞 인자 $\Theta(q)/\Theta^{(0)}(q)$ 가 $q'$ 에 무관이라 적분 \emph{밖}으로 빠진다:
\begin{equation*}
\int_0^q\mathcal K(q,q')\,b(q')\,\Theta(q')\,\dd q'\approx\frac{\Theta(q)}{\Theta^{(0)}(q)}\int_0^q\mathcal K(q,q')\,b(q')\,\Theta^{(0)}(q')\,\dd q'.
\end{equation*}
이를 $\Theta(q)=c(q)+(\text{위 적분})$ 에 넣고 $\Theta(q)$ 항을 좌변으로 모아 분모로 묶으면:
\begin{equation}
\boxed{\;\Theta_\mathrm A(q)=\frac{c(q)}{\,1-\dfrac{1}{\Theta^{(0)}(q)}\displaystyle\int_0^q\mathcal K(q,q')\,b(q')\,\Theta^{(0)}(q')\,\dd q'\,}\;.}
\label{eq:closed}
\end{equation}
\textbf{이것이 피팅에 쓰는 닫힌형이다} — 각 항의 의미: 분자 $c$ = 자유 감쇠($\Theta_\mathrm{init}$) + baseline
목표 기여, 분모의 적분 = 자기참조 감쇠($b<0$ 이라 분모 $>1$ → 꼬리를 줄임). \emph{무결성 점검}: 자기참조가
없으면($b\to0$) 분모 $\to1$, $\Theta_\mathrm A\to c$ 로 식~\eqref{eq:volterra} 의 올바른 극한(상수 목표서
$\Theta_0(1-e^{-q/L})$, 식~\eqref{eq:single_kernel})으로 복귀한다.
\begin{boundbox}
\textbf{Fredholm$\leftrightarrow$Volterra 가정차(\AL{16}; 그대로 쓰면 안 됨).} 사용자 논문(\emph{JCP}
\textbf{147}, 144111 (2017), Eq.(32)$\to$(34), \cite{jcp2017})은 \emph{Fredholm} 2종(고정 구간, 정상상태 확률)을
ratio-substitution 으로 닫는다. graphite 의 식~\eqref{eq:volterra} 는 \emph{Volterra}(인과 상한 $q$)다. 따라서
논문 결과를 그대로 옮기지 않고, 인과성을 유지한 채 자기참조만 선형화(단계 0)해 \emph{같은} ratio 구조를 적용했다.
또 \textbf{논문 자신이 ``reaction zone 이 넓으면 ratio 근사가 부정확''이라 경고}하는데, graphite 에서 넓은
spectrum $=$ 저온 stretched tail(가장 관심 영역)이므로 closure 가 가장 필요한 곳에서 가장 부정확할 수 있다.
$\Rightarrow$ \textbf{Plan A 단독 채택 금지}: Plan B(g-grid) 대비 상대오차
$\varepsilon=\lVert\Theta_\mathrm A-\Theta_\mathrm B\rVert/\lVert\Theta_\mathrm B\rVert$ 가 운용 $T\cdot$rate 에서
작을 때만 닫힌형으로 승격한다(BOUNDED).
\end{boundbox}

% ====================================================================
\section{ICA 와 DVA, 그리고 피팅 가능한 닫힌 논리식}\label{sec:fiteq}
% ====================================================================
\textbf{내부 전위 미분.} 식~\eqref{eq:charge_balance} 를 $q$ 로 미분하면(등온)
$Q_\cell=\dfrac{\partial Q_\bg}{\partial V_n}\dfrac{\dd V_n}{\dd q}+\sum_j Q_{j,\tot}\dfrac{\dd\xi_j}{\dd q}$,
따라서
\begin{equation}
\frac{\dd V_n}{\dd q}=\frac{Q_\cell-\sum_j Q_{j,\tot}\,\dd\xi_j/\dd q}{\partial Q_\bg/\partial V_n}.
\label{eq:dVdq}
\end{equation}
전체 phase progress 를 $Q_p\Theta\equiv\sum_j Q_{j,\tot}\xi_j$ 로 묶는다($Q_p=\sum_j Q_{j,\tot}$=용량 scale,
$\Theta\in[0,1]$ 무차원). 이 $\Theta$ 는 \S\ref{sec:kernel} 의 $\Theta$ 와 \emph{같은 양}이며, 그 post-peak
도함수가 $\dd\Theta_\mathrm{tail}/\dd q$(식~\eqref{eq:kernel_integral})다. 외부 전하 $\dd Q=\dd Q_\ext=
Q_\cell\,\dd q$ 는 배경 저장 $\dd Q_\bg=C_\bg\,\dd V_n$($C_\bg\equiv\partial Q_\bg/\partial V_n$)과
phase progress $Q_p\,\dd\Theta$ 의 합이다:
\begin{equation}
\dd Q=C_\bg\,\dd V_n+Q_p\,\dd\Theta.
\label{eq:dQsplit}
\end{equation}
이는 식~\eqref{eq:dVdq} 양변에 $C_\bg=\partial Q_\bg/\partial V_n$ 을 곱하고 $\dd q=\dd Q/Q_\cell$ 로 다시 쓴
\emph{동일식}이다(새 가정이 아니라 같은 전하보존 미분의 다른 표현). 양변을 $\dd Q$ 로 나누면 $1=C_\bg\,(\dd V_n/\dd Q)+Q_p\,(\dd\Theta/\dd Q)$ 이고, $\dd V_n/\dd Q$ 에 대해 풀어
역수를 취하면 내부 전위 기준 ICA
\begin{equation}
\boxed{\;\frac{\dd Q}{\dd V_n}=\frac{C_\bg(V_n,T)}{\,1-Q_p\,(\dd\Theta/\dd Q)\,}\;.}
\label{eq:fiteq}
\end{equation}
식~\eqref{eq:kernel_integral} 의 꼬리 기여가 $\dd\Theta/\dd Q$ 에 들어가 ICA 꼬리로 나타난다(좌표 bridge:
외부 전하 $Q=Q_\ext=Q_\cell q$ 이므로 $\dd\Theta/\dd Q=(1/Q_\cell)\,\dd\Theta/\dd q$). \emph{적용 영역}: post-peak
구간에선 꼬리 기여가 지배해 $\dd\Theta/\dd q\simeq\dd\Theta_\mathrm{tail}/\dd q$ 이고, peak 영역에선
$\dd\Theta/\dd q\to\dd\Theta_\eq/\dd q$(빠른 kinetics 극한, \S\ref{sec:kernel} 직전 single-mode 논의)로 환원된다.
$\dd\Theta/\dd Q$ 에 넣는 $\Theta$ 는 닫힌 $\Theta_\mathrm A$(식~\eqref{eq:closed}) 또는 Plan B 수치해를 $Q$ 로
미분한 값이다. 분모
$1-Q_p\,(\dd\Theta/\dd Q)>0$ 은 단조성(식~\eqref{eq:monotone}, $C_\bg=\partial Q_\bg/\partial V_n>0$)과
\emph{동치가 아니라 별개}의 유효범위 조건이다: 전자는 배경 용량의 국소 단조성(정적 chemical capacitance
양수)이고, 분모 양수성은 식~\eqref{eq:fiteq} 의 역수 $\dd V_n/\dd Q=(1-Q_p\,\dd\Theta/\dd Q)/C_\bg$ 가
양이라는 \emph{동역학 경로 조건}($C_\bg>0$ 전제하 $\dd V_n/\dd Q>0$)이다. 두 조건이 함께 성립해야
$\dd Q/\dd V_n>0$ 이다. staging peak 에서 분모가
$0$ 에 접근하면 $\dd Q/\dd V_n$ 이 발산하는데 이것이 정상적인 ICA peak 이며, 그 너머(분모 $<0$)는 전압 해의
특이점으로 보아 피팅에서 제외한다. 측정 축으로는
apparent 전위(식~\eqref{eq:Vapp})로 환산한다: 등온 정전류에서 $\dd V_{n,\app}/\dd q=\dd V_n/\dd q+
s_I|I|\,\partial R_n/\partial q$, 그리고 $\dd Q_\ext/\dd V_{n,\app}=Q_\cell/(\dd V_{n,\app}/\dd q)$,
$\dd V_{n,\app}/\dd Q_\ext$ 가 그 역수다(DVA). \emph{전압축 특성 꼬리 길이}는 다음과 같이 환산된다: post-peak 에서
$\dd V_n/\dd q$ 를 $q_a$ 둘레 상수로 근사하면(식~\eqref{eq:single_kernel} 유도의 가정 — 꼬리에서
$\dd\mathcal A_j/\dd q\simeq0$ 인 완만 영역) $V_n-V_n(q_a)\simeq(\dd V_n/\dd q)|_{q_a}(q-q_a)$ 이고,
single-mode 꼬리의 지수 $e^{-(q-q_a)/L}$ 가 $\exp[-(V_n-V_n(q_a))/L_\varphi]$ 로 옮겨져
\begin{equation}
L_\varphi\equiv\big|\dd V_n/\dd q\big|_{q_a}\,L
\label{eq:Lphi}
\end{equation}
가 정의된다($\dd V_n/\dd q$ 가 꼬리에서 완만하다는 가정 위의 BOUNDED 근사; 기호표에 $L_\varphi$[V] 등재).

\textbf{피팅 평가 순서(deliverable).} 식~\eqref{eq:fiteq} 는 다음 순서로 계산되는 닫힌 논리식이다:
\begin{enumerate}[label=\textbf{S\arabic*.},leftmargin=2.4em]
\item 전하 보존(식~\eqref{eq:charge_balance})으로 $V_n(q,T,\{\xi_j\})$ 를 구한다.
\item 평형 목표 $\xi_{\eq,j}(V_n,T)$(식~\eqref{eq:logistic}).
\item 속도 $k_j=k_0 e^{-(G-\chi_j\mathcal A_j)/RT}$, $\mathcal A_j=s_{\phi,j}F(V_{n,\drive}-U_j)$
  (식~\eqref{eq:kact}; 병목 시 식~\eqref{eq:kphys}).
\item 배리어$\to$길이 $L(G)$(식~\eqref{eq:L_of_G}) $\to$ spectrum $A_L$(식~\eqref{eq:spectrum}).
\item 꼬리 중첩 $\dd\Theta_\mathrm{tail}/\dd q$(식~\eqref{eq:kernel_integral}); 자기참조는 \emph{본 장에서} 닫는다 —
  Plan B(g-grid 수치, 항상) 또는 Plan A 닫힌형(식~\eqref{eq:closed}, $\varepsilon$ 통과 시).
\item $\dd Q/\dd V_n$(식~\eqref{eq:fiteq}) $\to$ 측정축 $\dd Q/\dd V_{n,\app}$.
\end{enumerate}
\textbf{피팅 파라미터}: $\{U_j,w_j,Q_{j,\tot},Q_\bg\}$(평형, 저속 OCV/GITT), $\{\Delta H_a,\Delta S_a,k_0\}$
(여러 $T$ 의 Arrhenius), $\chi_j$(여러 $|I|$/전위), $\rho_G$(독립 입력 — 역산 금지). \textbf{유효범위}:
$G-\chi_j\mathcal A_j\ge0$(Marcus, \AL{15}), $V_{n,\drive}\approx U_j+O(w_j)$. (전체 spectrum$\cdot$자기참조를
쓰는 정밀 평가의 solver 스킴$\cdot$수렴 판정은 부록 \S\ref{sec:ch1_numeric}.)

\subsection{심플 실데이터 피팅 근사식 (single-mode tail; 본 장만으로 사용 가능)}\label{sec:ch1_simplefit}
전체 spectrum 적분(식~\eqref{eq:kernel_integral}) 없이, \emph{하나의 우세 mode} 가 꼬리를 지배하는 극한
($A_L=\delta(L-L_0)$)에서는 위 모든 단계가 \textbf{종이와 연필로 닫히는} 가장 단순한 근사식으로 줄어든다. 이것이
실데이터 1차 피팅의 출발점이다. post-peak($q>q_a$) 에서 목표가 거의 정지하면(식~\eqref{eq:single_kernel}) 진행은
\begin{equation}
\Theta_\mathrm{tail}(q)\simeq\Theta_0\Big(1-e^{-(q-q_a)/L}\Big),\qquad
\frac{\dd\Theta_\mathrm{tail}}{\dd q}\simeq\frac{\Theta_0}{L}\,e^{-(q-q_a)/L},\qquad
\boxed{\;L=\frac{|I|}{Q_\cell\,k_j}=\frac{|I|}{Q_\cell k_0}\,e^{(G-\chi_j\mathcal A_j)/RT}\;}
\label{eq:simplefit}
\end{equation}
이고, 이를 fiteq(식~\eqref{eq:fiteq})에 넣으면 ICA 꼬리는 \emph{특성 길이 $L$ 의 지수 감쇠}로 나타난다.
\textbf{데이터 피팅 절차.}
\textbf{(0) 무대 고정(식별성 전제).} \S\ref{sec:falsify} 식별성 순서대로, 먼저 저속 OCV/GITT 로
$\{U_j,w_j,Q_{j,\tot},Q_\bg\}$ 를, 펄스/전류 차단으로 $R_n$ 을 고정한다($R_n$ 과 $\chi_j$ 가 공선형이므로 $R_n$ 을
\emph{먼저} 묶지 않으면 아래가 식별 불가). 측정 $V_{n,\app}$ 는 $V_n=V_{n,\app}-s_I|I|R_n$(식~\eqref{eq:Vapp} 역산,
$R_n$ 기지)으로 내부 $V_n$ 으로 \emph{직접} 환산한다(self-consistent 루프 아님).
\textbf{(1) $L$ 추출.} 꼬리 시작점 $q_a$ 를 \emph{먼저} 고정한다(ICA peak 위치, 또는 $\dd\xi_{\eq}/\dd q$ 가 임계
이하로 떨어지는 첫 $q$) — $q_a$ 를 자유 파라미터로 두면 $L$ 과 공선형이다. 그 뒤 post-peak ICA 꼬리를
모델 $y(q)=A\,e^{-(q-q_a)/L}+(\text{baseline})$ 로 회귀해 \emph{한 숫자} $L$ 을 뽑는다 — 진폭 $A$ 는 nuisance
($\Theta_0\cdot Q_p/C_\bg$ 등을 흡수; 관심량은 $L$), 단일-mode 꼬리에선 fiteq 분모 $1-Q_p\,\dd\Theta/\dd Q\!\approx\!1$
이라 ICA 꼬리 $\approx(Q_p/C_\bg)(\Theta_0/L)e^{-(q-q_a)/L}$ 로 닫힌다(전압축이면 $L_\varphi$, 식~\eqref{eq:Lphi}).
\textbf{(2) 유효 엔탈피(prefactor$\cdot U_j(T)$ 보정).} $1/L\propto k_j=k_0(T)e^{-(G-\chi_j\mathcal A_j)/RT}$ 이고
$k_0=(k_BT/h)\kappa$(Eyring, 식~\eqref{eq:kact})에 \emph{T-선형 prefactor} 가 있어, $\ln(1/L)$ 를 그대로
$1/T$ 로 회귀하면 prefactor 의 $\ln T$ 항이 기울기를 오염시킨다(순수 $\Delta H_a$ 가 아니다). 표준 Eyring 선형화
($\ln(k/T)$ 대 $1/T$)대로 prefactor 를 보정하고, 또한 $\mathcal A_j=s_{\phi,j}F(V_{n,\drive}-U_j(T))$ 의 $U_j(T)$
T-의존 때문에 $\chi_j\mathcal A_j(T)/RT$ 가 $1/T$ 에 \emph{비선형}으로 섞이므로 이를 점별 기지값으로 먼저 빼낸다. 즉
\begin{equation*}
y(T)\equiv\ln\!\frac{1}{L\,T}-\frac{\chi_j\,\mathcal A_j(T)}{RT}\quad\text{를}\quad \frac1T\ \text{로 회귀} \;\Rightarrow\; \text{기울기}=-\frac{\Delta H_a}{R}\ (\text{순수}).
\end{equation*}
$\chi_j\mathcal A_j(T)$ 는 각 $T\cdot$측정점에서 (0)$\cdot$(3)단계 기지값으로 계산한다($\chi_j$ 는 (3)에서 먼저 얻으므로
비순환). prefactor 의 $\kappa(T)$ 약한 T-의존을 무시하는 근사는 $\Delta H_a/(RT)\gg1$($\Delta H_a\gtrsim40$ kJ/mol)
에서 BOUNDED 다. (보정 전 기울기 $-(\Delta H_a-\chi_j\mathcal A_j)/R$ 는 \emph{유효} 엔탈피로만 해석한다.)
(부록 \S\ref{sec:ch6_arrhenius} 식~\eqref{eq:ch6_arrhenius} 의 raw $\ln(1/L)$ 분해에 들어 있는 $+\chi_j\mathcal A_j/RT$ 를
여기서 빼 순수화한 것이며, 보정 전 유효 기울기 $-(\Delta H_a-\chi_j\mathcal A_j)/R$ 는 식~\eqref{eq:ch6_arrhenius_slope} 와 동일.)
\textbf{(3) $\chi_j$.} \emph{peak 근방}($\mathcal A_j=s_{\phi,j}F(V_{n,\drive}-U_j)$ 가 유한한 영역)에서 내부 전위
$V_{n,\drive}=V_n$(0단계 환산)를 가로축으로 $\partial\ln L/\partial V_{n,\drive}=-\chi_j s_{\phi,j}F/RT$ 기울기에서
$\chi_j$ 를 얻는다. \textbf{비순환 보장}: $\mathcal A_j=s_{\phi,j}F(V_{n,\drive}-U_j)$ 는 각 데이터점에서 (0)단계 환산
$V_n$ 과 (0)단계 고정 $U_j(T)$ 로 \emph{계산되는 기지값}(자유 파라미터 아님)이다. 따라서 순서는 비순환이다 — (3)에서
$\chi_j$ 를 먼저 얻고, 그 $\chi_j$ 로 (2)의 종속변수 $y(T)$ 에서 점별 보정항 $\chi_j\mathcal A_j(T)/RT$ 를 \emph{빼서} 순수
$\Delta H_a$ 를 분리한다((2)↔(3) 닭-달걀 아님). \textbf{즉 $\{L,\Delta H_a,\chi_j\}$ 가 단일 지수 꼬리 피팅
+ 식별성 순서로 나온다}(필요한 근거가 모두 본 장 안에 있다: 근사식~\eqref{eq:simplefit}, 식별성 \S\ref{sec:falsify}). 꼬리가 단일 지수에서
벗어나(stretched) 휘면 그때 비로소 spectrum(식~\eqref{eq:spectrum})$\cdot$자기참조(식~\eqref{eq:closed})가
필요하다 — \emph{단일 지수 근사의 실패 자체가 분포$\cdot$결합의 신호}다(반증, \S\ref{sec:falsify}).

\textbf{C-rate 기준식(0.2C anchor).} 전류 의존 분극을 저율 기준 곡선에 상대화하면, $0.2$C($I_{0.2C}>0$)를 anchor 로
\begin{equation}
V_{n,\app}(q;|I|)\approx V_{n,0.2C}(q)+s_I\big[\,|I|\,R_n(q,T,|I|)-I_{0.2C}R_n(q,T,I_{0.2C})\,\big],
\label{eq:02c}
\end{equation}
전류 의존이 약하면 단순형 $V_{n,\app}\approx V_{n,0.2C}+s_I(|I|-I_{0.2C})R_n(q,T)$ 다 — $0.2$C 에서 상변이 지연이
충분히 작거나 그 지연을 기준 곡선에 흡수해도 되는 경우의 실용 근사. 여기에 식~\eqref{eq:simplefit} 의 $L(|I|)$ 단축
($\partial\ln L/\partial V_{n,\drive}<0$)을 더하면 C-rate 가 커질수록 꼬리가 짧아지고 peak 겹침이 심해지는 관측이
정량 설명된다. \emph{피팅 절차와의 연결}: 식~\eqref{eq:02c} 는 여러 $|I|$ 의 $V_{n,\drive}$ 를 공통 저율 기준으로
정렬해 위 (3)단계 $\chi_j$ 회귀의 가로축 입력을 일관되게 만드는 데 쓴다.

\textbf{초기값(EMG 등 경험식).} 비선형 피팅의 \emph{초기값}으로 $\dd Q/\dd V=B(V)+\sum_j A_j\,g_\mathrm{EMG}
(V;\mu_j,\sigma_j,\lambda_j)$ 같은 peak 함수를 쓸 수 있으나 주 모델이 아니다 — 꼬리는 $k_j,\rho_G,R_n$$\cdot$전하
보존의 결합 결과이며, EMG 의 $\lambda_j$(지수 꼬리율)는 식~\eqref{eq:simplefit} 의 $1/L$ 의 경험적 대응물로만 본다.

% ====================================================================
\section{온도 의존 꼬리의 분해와 반증(falsification)}\label{sec:falsify}
% ====================================================================
온도에 따른 꼬리 증감을 한 항으로 단정하지 않고 \textbf{세 계층}의 결합으로 본다.
\begin{longtable}{p{0.24\textwidth} p{0.30\textwidth} p{0.36\textwidth}}
\toprule 계층 & 담당 항 & 해석 \\ \midrule \endhead
평형 열역학 & $U_j(T),w_j(T),Q_\bg(V_n,T)$ & 전이 중심 전위$\cdot$평형 폭$\cdot$배경 용량의 온도 의존 \\
동역학 지연 & $k_j(T,I),\rho_G(g;T)$ & 진행률이 평형을 못 따라가 생기는 꼬리/broadening \\
전압축 분극 & $R_n(q,T,|I|)$ & apparent 전위 이동, 기울기 변화, 분극 peak 변형 \\
\bottomrule
\end{longtable}
\textbf{전위 의존 spectrum shift.} 식~\eqref{eq:L_of_G} 에서 $\partial\ln L/\partial V_{n,\drive}
=-\chi_j s_{\phi,j}F/RT<0$(방전 방향, $\chi_j\ge0$): 구동 전위가 커지면 $L$ 이 작아져 꼬리가 짧아진다 — C-rate/전위가
꼬리를 어떻게 바꾸는지의 정량 예측이다.

\textbf{반증 규칙(forward-prediction 검정; 역산 금지).} 본 이론은 다음으로 \emph{반증 가능}하다.
\begin{enumerate}[label=(N\arabic*),leftmargin=2.4em]
\item \textbf{평형 형태}: 저속 OCV 의 near-peak 감쇠가 $\log(dQ/dV)$ 대 $(V-U)$ 면 logistic, 대 $(V-U)^2$
  면 가우시안 — 어느 쪽도 안 맞으면 평형 ansatz 재검토.
\item \textbf{Arrhenius}: $\ln k_j$(또는 $\ln(1/L)$)대 $1/T$ 가 직선이 아니면 단일 활성화 그림 기각.
\item \textbf{전위 의존}: $\partial\ln L/\partial V_{n,\drive}$ 가 위 예측과 다르면 $\chi_j$ 해석 기각.
\item \textbf{비유일성 차단}: $\rho_G$ 를 꼬리에서 역산하지 않고 독립 입력으로만 — 역산하면 비유일이라
  반증력이 없다.
\item \textbf{비퇴화 — transport$\cdot$charge-transfer 와의 분리(★ 본 이론의 가장 약한 지점).} 꼬리 길이 scaling
  $L\propto|I|$ 과 ``전위 커지면 꼬리 짧아짐''은 \emph{단순 transport 분극$\cdot$charge-transfer 도 공유}하므로,
  그것만으로 barrier-spectrum 에 귀속하면 \emph{퇴화}(오귀속)다. 구별 signature(두 조건 (a)$\cdot$(b)가 \emph{모두}
  필요하며, $k_{j,\lim}$ 이 전위무관일 때만 성립): (a) \emph{전위 의존 Arrhenius slope} — 보정 전 측정 기울기
  $\Delta H_a^{\eff}=\Delta H_a-\chi_j\mathcal A_j$(식~\eqref{eq:simplefit}; 추출은 \S\ref{sec:ch1_simplefit} 절차 (2))
  가 구동 전위에 따라 이동함. \emph{단 이것만으로는 transport 와 여전히 퇴화할 수 있다}: 고체상 확산이 율속이고
  확산계수가 점유율 의존($D=D(\xi)$, thermodynamic factor 경유)이면 transport-limited 유효 엔탈피도 전위에 따라
  이동한다(식~\eqref{eq:kphys} 의 $k_{j,\lim}$ 전위의존). 따라서 (a)는 GITT relaxation transient 의 형상
  ($\sqrt t$ Cottrell=확산 율속 vs 지수=계면 율속)으로 $k_{j,\lim}$ 의 전위의존을 독립 판정해 \emph{배제}한 뒤에만
  barrier-lowering signature 로 유효하다; (b) charge-transfer 과전압 $\eta_{\mathrm{ct}}$ 도 전위 지수의존이라 단일
  전극서 $\chi_j$ 와 공선형이므로 — \textbf{GITT/전류 차단으로 $\eta_{\mathrm{ct}}(R_{\mathrm{ct}})$ 를 독립 분리한
  뒤에만} $\chi_j$ 의 ``배리어 낮춤'' 귀속이 성립한다(단순 rate-broadening \cite{dubarry2022,fly2020} 과의 구별이
  핵심). 이 분리(확산 율속 배제 + $\eta_{\mathrm{ct}}$ 분리) 없이 단일 데이터셋에서 꼬리를 barrier-spectrum 으로
  귀속하는 것은 금지한다.
\end{enumerate}
\textbf{파라미터 제한 순서(식별성).} 모든 파라미터를 한 데이터에서 동시에 자유 피팅하면 공선형(co-linear)이라
분리 불가다. 다음 \emph{순차} 제약으로 식별성을 확보한다: \textbf{(1)} 저속 OCV/GITT(준평형, $|I|\to0$)로
$\{U_j,w_j,Q_{j,\tot},Q_\bg\}$ 를 먼저 고정 — 동역학 꼬리가 사라진 평형 무대를 잡는다. \textbf{(2)} 펄스/전류
차단으로 $R_n$(전압축 분극)을 독립 제한 — $R_n$ 과 $\chi_j$ 는 둘 다 꼬리를 늘려 공선형이므로 $R_n$ 을 먼저
묶는다. \textbf{(3)} 여러 온도의 Arrhenius($\ln(1/L)$ 대 $1/T$)로 $\{\Delta H_a,\Delta S_a,k_0\}$. \textbf{(4)}
여러 전류/전위($\partial\ln L/\partial V_{n,\drive}$)로 $\chi_j$. \textbf{(5)} $\rho_G$ 는 독립 입력(역산 금지,
N4). 이 순서를 어기면 식별성이 무너져 \emph{어떤 적합도라도} 물리 해석이 불가능하다(과적합).

% ====================================================================
\section{부록: 자기참조식의 수치 평가와 검증 (Plan B 상세)}\label{sec:ch1_numeric}
% ====================================================================
전체 spectrum$\cdot$자기참조를 쓰는 정밀 평가는 Plan B(직접 $g$-grid)로 푼다. 단순 꼬리는 본문
\S\ref{sec:ch1_simplefit} 의 단일 지수 근사로 충분하나, stretched 꼬리(저온)나 강한 자기참조에서는 본 수치
평가가 기준해다.
\textbf{스킴.} (i) 배리어 격자 $\{g_m\}_{m=1}^{N_g}$ 를 관심 $T\cdot$전위에서 $L(g_m)$(식~\eqref{eq:L_of_G})이
관측 윈도를 덮도록 — 특히 large-$L$(긴 꼬리) 영역 — 잡는다. (ii) mode 별 ODE
$\dd\xi_j(g_m)/\dd t=k_j(g_m)[\xi_{\eq,j}-\xi_j(g_m)]$ 를 시간 적분하되, 매 시점 전하 보존
(식~\eqref{eq:charge_balance})의 대수 제약으로 $V_n$ 을 root-find 한다 — 미분변수 $\xi_j(g_m)$ + 대수 제약
$V_n$ 의 \emph{index-1 DAE} 이며, 제약을 한 번 미분하면 $\partial Q_\bg/\partial V_n\ge\epsilon_Q>0$
(식~\eqref{eq:monotone})이 $\dd V_n/\dd t$ 를 유일하게 주므로 index 가 정확히 1 이다. (iii)
$\xi_j(t)=\sum_m w_m\rho_G(g_m)\xi_j(g_m,t)$ 로 평균.
\textbf{근거 기반 tolerance(임의 상수 아님).} root-find 종료는 전하 보존 잔차 $<\delta_q Q_\cell$($\delta_q$=좌표
분해능), 시간적분은 $\xi\in[0,1]$ 상대오차와 측정 ICA 잡음 floor 에서 잡는다 — cap/clip 이 아닌 측정$\cdot$물리
scale 기준.
\textbf{Plan A 승격 검증.} 닫힌형 $\Theta_\mathrm A$(식~\eqref{eq:closed})는 Plan B 대비
$\varepsilon=\lVert\Theta_\mathrm A-\Theta_\mathrm B\rVert/\lVert\Theta_\mathrm B\rVert$ 를 운용 $T\cdot$rate$\cdot$
tail-window 에서 측정해 $\varepsilon\le\varepsilon_{\mathrm{tol}}=\max(\varepsilon_{\mathrm{disc}},\varepsilon_{\mathrm{noise}})$
(기준해 이산화 오차 $\vee$ 측정 잡음)일 때만 승격한다. \emph{특정 수치값}은 데이터 의존이라 미리 고정하지 않는다(FLAGGED). 넓은
spectrum(저온 stretched)에서 ratio 근사가 가장 열화하므로(\cite{jcp2017} 자기명시) 이 검증이 필수다.
\textbf{격자 수렴.} $\lVert\Theta_{N_g}-\Theta_{2N_g}\rVert/\lVert\Theta_{2N_g}\rVert\le\varepsilon_{\mathrm{disc}}$
(Plan A 승격에 쓰는 이산화 오차 기준 $\varepsilon_{\mathrm{disc}}$ 와 동일값)가 될 때까지 $N_g$ 를 배가한다 —
정성적 ``수렴할 때까지''가 아니라 위 $\varepsilon_{\mathrm{tol}}$ 체인과 닫힌 정량 종료조건이다. solver 코드 구현
자체는 본 이론 범위 밖(P1; 이론식$\cdot$검증 원리만 제시).

% ====================================================================
% 부록 B — 유도식의 실데이터 피팅 워크플로 (구 Chapter 6 흡수: Ch1 자기완결)
% ====================================================================
\section*{\normalfont\Large 부록 B — 유도식의 실데이터 피팅 워크플로 (구 Chapter 6)}
\label{sec:ch1_fitting_practice}
\addcontentsline{toc}{section}{부록 B — 실데이터 피팅 워크플로}
\noindent\emph{아래 절들은 구 Chapter 6 ``피팅 실무 부록''을 본 장에 흡수한 것이다(Ch.6 해체). 본문 \S1--\S15 의 이론과 심플 근사식(\S\ref{sec:ch1_simplefit}) 위에서, 실데이터 피팅을 운영하는 식별성$\cdot$수치 해법(index-1 DAE)$\cdot$순차 제약(저속 OCV$\to$GITT$\to$다온도 Arrhenius$\to$C-rate)$\cdot$반증 워크플로를 담는다. 가정 원장 AL-60$\sim$69 를 포함한다. 발열$\cdot$히스테리시스 피팅은 각각 Ch.2$\cdot$Ch.4, Ch.5 의 정의식을 참조한다.}
\section{부록 B 도입 — 본문 식과 피팅 파라미터 정리}\label{sec:ch6_inherit}
% ====================================================================
본 장은 Chapter 1--5(RB 재구성본)의 다음을 \textbf{수정 없이} 입력으로 받는다(프로젝트 검수 P3-6, P5). 각 식
번호는 앞 장의 boxed 결과를 가리키며, 본 장은 그 위에 \emph{피팅 방법론} 계층을 적층한다(앞 장 본문 기본식을
고쳐 쓰지 않는다).
\begin{linkbox}
\textbf{(L1) 전하 보존식(중심 상태식, Ch1).} $Q_\cell\,q=Q_\bg(V_n,T)+\sum_{j=1}^{N_p}Q_{j,\tot}\xi_j$. $V_n$ 은
그 해(외부 입력 아님); 단조성 $\partial Q_\bg/\partial V_n\ge\epsilon_Q>0$.\\
\textbf{(L2) 진행률 동역학(Ch1$\cdot$Ch3).} $\dfrac{\dd\xi_j}{\dd t}=k_j[\xi_{\eq,j}(V_n,T)-\xi_j]$ 또는 forward/
backward 형 $r_j^+(1-\xi_j)-r_j^-\xi_j$; detailed balance $r_j^+/r_j^-=\xi_{\eq,j}/(1-\xi_{\eq,j})$;
$k_j=k_0(T)e^{-(G-\chi_j\mathcal A_j)/RT}$, $\mathcal A_j=s_{\phi,j}n_j^{\eff}F(V_{n,\drive}-U_j)$, 기본 $V_{n,\drive}=V_n$.\\
\textbf{(L3) 피팅 논리식과 평가 순서(Ch1).} $\dfrac{\dd Q}{\dd V_n}=\dfrac{C_\bg(V_n,T)}{1-Q_p\,(\dd\Theta_\mathrm{tail}/\dd Q)}$,
$C_\bg=\partial Q_\bg/\partial V_n$; 평가 순서 S1--S6 (charge-balance root $\to\xi_{\eq}\to k_j\to L(G)\to A_L\to$
kernel integral $\to\dd Q/\dd V_n\to$ 측정축 $\dd Q/\dd V_{n,\app}$).\\
\textbf{(L4) self-consistent 적분식(Ch1).} $\Theta(q)=\Theta_\mathrm{init}(q)+\int_0^q\mathcal K(q,q')\,
\Theta_\eq(V_n(q',\Theta(q')),T)\,\dd q'$, $\mathcal K(q,q')=\int_0^\infty A_L(L)\tfrac1L e^{-(q-q')/L}\dd L$ (Volterra 2종).\\
\textbf{(L5) closure 2-tier(Ch1).} \emph{직접 수치적분}(분포 $g$-grid)$=$항상 보존되는 기준 해법; \emph{Plan A}
$=$Fredholm ratio-substitution closed-form(사용자 PhD, \cite{lee2011,son2013,jcp2017})$=$검증 통과 시 가속 후보.\\
\textbf{(L6) 발열(Ch2$\cdot$Ch4).} $\dot{\mathcal Q}_\irr=\sum_j n_j^{\eff}I_j\eta_j\ge0$, $\dot{\mathcal Q}_\rev=
s^\conv|I|T\,\partial U/\partial T$ 류; calorimetry 부재 시 \textbf{forward prediction} 으로만(피팅 금지).\\
\textbf{(L7) 히스테리시스(Ch5).} signed current $I_s$$\cdot$$q_\stt$, 충전 부호 $s_{\phi,j}^b=-1$, branch target
$\xi_{\tar,j}^b=\xi_{\sstat,j}^b$, $\Delta V_\obs\ne\Delta V_\hys$, loop-area 분리 발열.\\
\textbf{(L8) 반증(Ch1).} N1 Arrhenius 직선성, N2 전위 의존 spectrum shift($\partial\ln L/\partial V_{n,\drive}<0$),
N3 저온 lineshape(평형 broadening vs disorder), N4 $\rho_G$ forward-only(역산 금지).
\end{linkbox}
\noindent\textbf{$\conv,\stt,\hys,\OCV$ 등 약어는 Ch1--5 정의 그대로}이며 본 장은 이를 재정의하지 않는다.
Chapter 1--5 본문 기본식은 수정하지 않는다(P3-6). 본 장은 이들을 \emph{데이터로 제약}하는 절차로 배치한다.

\textbf{피팅 파라미터의 전수 목록(앞 장이 남긴 것).} 본 장이 식별 대상으로 삼는 파라미터는 다음으로 닫힌다(이 외의
자유 knob 추가 금지). 각각 \emph{어느 장이 정의했고 어느 독립 실험이 1차로 제약하는지}를 \S\ref{sec:ch6_identif}
에서 매핑한다.
\begin{longtable}{p{0.20\textwidth} p{0.30\textwidth} p{0.44\textwidth}}
\toprule 파라미터 & 정의 장$\cdot$식 & 물리적 의미(피팅 시 해석) \\ \midrule \endhead
$U_j,\ w_j$ & Ch1 eq:logistic & 전이 중심 전위$\cdot$평형 전이 폭(열역학 무대) \\
$Q_{j,\tot},\ Q_\bg(V_n,T)$ & Ch1 eq:charge\_balance & 전이 용량 기여$\cdot$배경 누적 용량(전하 보존) \\
$\Delta H_{a,j},\ \Delta S_{a,j},\ k_0$ & Ch1 eq:kact (\AL{10}) & 활성화 엔탈피$\cdot$엔트로피$\cdot$prefactor(Eyring) \\
$\chi_j(\equiv\beta_j$; 조건부, 대칭 Marcus$\cdot$정상 target 극한 한정$)$ & Ch1/Ch3 (\AL{11},31) & 배리어 낮춤 민감도 $=$ transfer coefficient \\
$\rho_G(g;T)$ & Ch1 eq:spectrum (\AL{13}) & 배리어 확률밀도(\textbf{독립 입력}; 역산 금지, \AL{16}) \\
$R_n$ & Ch1 eq:Vapp / Ch3 & 음극 apparent 분극 계수($\ne R_\ct\ne R_{n,\eff}$) \\
$k_{j,\lim}$ & Ch1 eq:kphys & 수송$\cdot$확산$\cdot$유한자리 병목(독립 제약 시만) \\
$n_j^{\eff}=RT/(Fw_j)$ & Ch3 (\AL{34}) & effective 전자수(detailed balance 가 강제; 독립 자유 X) \\
$s_{\phi,j}^b,\ \xi_{\tar,j}^b,\ h_j$ & Ch5 (\AL{53}--56) & branch 부호$\cdot$target$\cdot$hysteresis state(충전 시) \\
\bottomrule
\end{longtable}

% ====================================================================
\section{보유 데이터와 식별 가능 / 불가 항목의 분리}\label{sec:ch6_data}
% ====================================================================
\textbf{피팅의 첫 단계는 회귀가 아니라 \emph{무엇을 제약할 수 있는지}의 분리}다. 1차 검증 기준은 사용자 보유
\textbf{음극 반쪽셀 GITT, STO--OCV 테이블, $0.1/0.2/0.5/1.0\,\mathrm{C}$ rate series} 다(\AL{63}). 이 데이터는
모델을 \emph{전부} 결정하지 않으므로, 제약 가능/불가 항을 먼저 나눠 \emph{과결정(over-fitting)} 을 차단한다.
\begin{longtable}{p{0.28\textwidth} p{0.34\textwidth} p{0.32\textwidth}}
\toprule 데이터 & 우선 제약 가능 & 직접 확정 곤란 \\ \midrule \endhead
반쪽셀 STO--OCV & $V_{n,\OCV}(q),U_j,w_j,Q_{j,\tot},Q_\bg$ 열역학 골격 & $k_j,\rho_G(g)$, transport limitation \\
반쪽셀 GITT & instantaneous polarization($R_n$ 일부), relaxation signature, OCV 정합 & $k_j$ 단독값, solid diffusion vs phase relaxation 완전 분리 \\
$0.1/0.2\,\mathrm C$ & 저율 quasi-equilibrium proxy, peak 위치/폭 & strict thermodynamic equilibrium \\
$0.5/1.0\,\mathrm C$ & kinetic lag, broadening, transport stress test & pure frozen kinetic limit \\
다온도 series & Arrhenius 기울기($E_a$), $\Delta S_a$ & 단일 온도만으로는 $T$-의존 분리 불가 \\
온도별 수명/full-cell & 정성적 열화 trend & $U_j(T),w_j(T),k_j(T),\Delta S_j$ 정량 분리 \\
\bottomrule
\end{longtable}
\begin{boundbox}
\textbf{데이터 제약(\AL{63}).} 단일 온도 자료로 $U_j(T),w_j(T),k_j(T)$ 의 온도 의존을 분리하지 않는다 — Arrhenius
기울기는 \emph{여러} 온도가 있어야 정의된다(\S\ref{sec:ch6_arrhenius}). calorimetry 없이 Ch2/Ch4 발열 파라미터를
피팅하지 않는다 — 발열 항은 \textbf{forward prediction}(미시$=$거시 $n^{\eff}I_j\eta_j$ 로 예측만)으로 두고
calorimetry 확보 후 검증한다(\AL{67}).
\end{boundbox}
\begin{boundbox}
\textbf{등온 가정 유효 범위.} 등온 피팅은 온도 상승이 충분히 작거나 외부 온도 제어가 강할 때만 유효하다. 고율
자기발열로 $T(t)$ 가 유의하게 변하면 $U_j(T),w_j(T),k_j(T),R_n(T),k_{j,\lim}(T),\rho_G(g;T)$ 가 함께 변하므로
등온 C-rate 자료만으로 kinetic parameter 를 확정하지 않고 Ch2/Ch4 thermal coupling 으로 돌아간다(\AL{67}).
\end{boundbox}

% ====================================================================
\section{피팅 평가 순서 S1--S6 의 구현 (Ch1 평가 순서의 피팅 워크플로화)}\label{sec:ch6_S}
% ====================================================================
\textbf{물리.} Ch.1 의 결론식(eq:fiteq)은 평가 순서 S1--S6 으로 \emph{한 점}의 $\dd Q/\dd V_n$ 을 계산한다. 피팅은
이 계산을 파라미터의 함수로 두고, 측정 ICA 와 비교해 residual 을 최소화한다. 본 절은 S1--S6 각 단계가 어떤
\emph{수치 연산}으로 실현되는지를 명시한다 — 이 수치 연산이 곧 피팅의 inner loop(주어진 파라미터에서 모델
ICA 를 한 번 평가하는 forward 계산)이다.

\subsection{S1 — 전하 보존식 root 로 $V_n$ 결정 (피팅 inner-loop 의 첫 연산)}\label{sec:ch6_S1}
주어진 진행 좌표 $q$ 와 진행률 $\bm\xi=(\xi_1,\dots,\xi_{N_p})$ 에서 내부 전위 $V_n$ 은 Ch.1 전하 보존식을
$V_n$ 에 대해 푸는 \emph{root finding} 으로 얻는다(외부에서 읽는 값이 아니다, Ch1):
\begin{equation}
F_\cb^\dyn(V_n;q,\bm\xi)\;=\;Q_\bg(V_n,T)+\sum_{j=1}^{N_p}Q_{j,\tot}\xi_j-Q_\cell q\;=\;0.
\label{eq:ch6_dyn_root}
\end{equation}
$\bm\xi$ 는 이 root finding 에서 \emph{독립 입력}(시간 적분이 따로 갱신)이므로 Jacobian 은 배경 용량 기울기뿐이다:
\begin{equation}
\frac{\partial F_\cb^\dyn}{\partial V_n}=\frac{\partial Q_\bg}{\partial V_n}\;\ge\;\epsilon_Q>0.
\label{eq:ch6_dyn_deriv}
\end{equation}
Ch.1 의 단조성 floor(eq:monotone)가 유지되면 해 존재 범위에서 root 는 유일하다(중간값 정리 $+$ 단조성, Ch1).

\textbf{평형 OCV root 와의 구분(피팅에서 중요).} 평형/준평형에서는 $\xi_j=\xi_{\eq,j}(V_n)$ 이므로 root 식이
달라지고 Jacobian 에 $\xi_{\eq}$ 항이 추가된다:
\begin{equation}
F_\cb^\OCV(V_n;q)=Q_\bg(V_n)+\sum_j Q_{j,\tot}\xi_{\eq,j}(V_n)-Q_\cell q=0,\qquad
\frac{\partial F_\cb^\OCV}{\partial V_n}=\frac{\partial Q_\bg}{\partial V_n}+\sum_j Q_{j,\tot}\frac{\partial\xi_{\eq,j}}{\partial V_n}.
\label{eq:ch6_ocv_root}
\end{equation}
\textbf{두 root 의 derivative 가 다르므로(식~\eqref{eq:ch6_dyn_deriv} vs \eqref{eq:ch6_ocv_root}) 같은 Newton
Jacobian 으로 처리하면 안 된다.} 식~\eqref{eq:ch6_ocv_root} 의 $\partial F_\cb^\OCV/\partial V_n$ 은 Ch.2 의
평형 chemical capacitance(OCV 온도계수 유도에 쓴 $\dd Q/\dd V_{n,\OCV}$)와 같은 양이다 — 즉 \textbf{S1 의
평형 극한은 Ch.2 의 평형 미분용량과 자동 정합}하며, 둘이 어긋나면 피팅이 잘못된 것이다(내부 정합 점검).
\begin{verifybox}
\textbf{root-find 종료 기준은 임의 상수가 아니다(\AL{64}, GS-4$'$).} Newton 종료는 ``residual 이 충분히 작다''가
아니라 \emph{전하 단위의 물리 scale} 로 정규화한다: $|F_\cb^\dyn|\le\varepsilon_Q^{\mathrm{root}}\equiv
\delta_q\,Q_\cell$, 여기서 $\delta_q$ 는 진행 좌표 $q$ 의 측정/계산 분해능(예: $\dd q$ 격자 간격)이다. 즉 ``배경
용량 잔차가 한 격자 스텝의 외부 전하보다 작다''는 \emph{측정 분해능 기준}이며, 무근거 숫자가 아니다. 부동소수점
한계상 $\varepsilon_Q^{\mathrm{root}}$ 을 기계 epsilon 아래로 두지 않는다($\sqrt{\varepsilon_{\mathrm{mach}}}$
근방이 Newton 의 통상 floor, \cite{brenan1996}).
\end{verifybox}

\subsection{S2 — 평형 target $\xi_{\eq,j}$ (logistic 평가)}\label{sec:ch6_S2}
S1 의 $V_n$ 에서 Ch.1 logistic(eq:logistic) $\xi_{\eq,j}(V_n,T)=[1+\exp(-s_{\phi,j}(V_n-U_j)/w_j)]^{-1}$ 를 평가한다.
피팅에서 $U_j,w_j$ 는 S1 의 $Q_{j,\tot},Q_\bg$ 와 함께 \emph{저속 OCV} 가 1차로 제약한다(\S\ref{sec:ch6_ocvfit}).
$n_j^{\eff}=RT/(Fw_j)$ 는 $w_j$ 가 정해지면 \emph{종속}이며 독립 자유 파라미터가 아니다(Ch3 \AL{34}).

\subsection{S3 — 속도상수 $k_j$ (Eyring 평가, 병목 직렬 합성)}\label{sec:ch6_S3}
$\mathcal A_j=s_{\phi,j}n_j^{\eff}F(V_{n,\drive}-U_j)$(기본 $V_{n,\drive}=V_n$), $k_j^{\mathrm{act}}=k_0(T)
e^{-(G-\chi_j\mathcal A_j)/RT}$ 를 평가한다(Ch1 eq:kact). 병목이 유효하면 Ch.1 의 \emph{직렬 시간척도 합}
$1/k_j^{\eff}=1/k_j^{\mathrm{act}}+1/k_{j,\lim}$ 로 합성한다(eq:kphys). \textbf{강구동에서 $k_j^{\mathrm{act}}$ 가
커지는 것을 임의 cap 으로 자르지 않는다}(GS-4); 수치 stiffness 는 \S\ref{sec:ch6_dae} 의 implicit solver 가
흡수한다.

\subsection{S4 — 배리어$\to$길이 $L(G)$ 와 spectrum $A_L$}\label{sec:ch6_S4}
Ch.1 의 닫힌 지수 매핑(eq:L\_of\_G) $L(G)=\dfrac{|I|}{Q_\cell k_0(T)}\exp[(G-\chi_j\mathcal A_j)/RT]$ 과 변수변환
밀도 보존(eq:spectrum) $A_L^{\mathrm{prob}}(L)=\rho_G(G(L))\,(RT/L)\,H(L-L_{\min})$ 를 평가한다. \textbf{Heaviside
$H$ 는 Ch.1 변수변환의 자코비안 support 하한($G\ge0\Leftrightarrow L\ge L_{\min}$)에서 \emph{유도}된 정의역
경계}이며, 피팅 편의를 위한 임의 절단이 아니다(Ch1 \S spectrum; GS-4 비위반). 수치적으로는 $g$-grid 의 하한을
$L_{\min}$ 에 맞추어 음의 배리어 영역을 적분에서 제외한다.

\subsection{S5 — 꼬리 중첩 $\dd\Theta_\mathrm{tail}/\dd q$ (kernel integral; 자기참조 풀이는 \S\ref{sec:ch6_grid})}\label{sec:ch6_S5}
Ch.1 의 kernel integral(eq:kernel\_integral) $\dfrac{\dd\Theta_\mathrm{tail}}{\dd q}=\int_0^\infty A_L(L)\,\tfrac1L\,
e^{-(q-q_a)/L}\,\dd L$ 를 평가한다. $\xi_j$ 와 $V_n$ 이 전하 보존으로 얽혀 있어 이는 self-consistent Volterra
2종(L4)이며, 그 \emph{직접 수치 풀이가 본 장 closure 의 기준 해법}이다(\S\ref{sec:ch6_grid}--\ref{sec:ch6_closure}).

\subsection{S6 — $\dd Q/\dd V_n$ 과 측정축 $\dd Q/\dd V_{n,\app}$}\label{sec:ch6_S6}
Ch.1 결론식(eq:fiteq) $\dd Q/\dd V_n=C_\bg/(1-Q_p\,\dd\Theta/\dd Q)$ 를 평가하고, 측정축으로는 apparent
전위(eq:Vapp) $V_{n,\app}=V_n+s_I|I|R_n$ 로 환산한다: 등온 정전류에서 $\dd V_{n,\app}/\dd q=\dd V_n/\dd q+
s_I|I|\,\partial R_n/\partial q$, DVA 는 그 역수다. \textbf{측정 ICA 는 $\dd Q/\dd V_{n,\app}$ 와 비교}하며, 모델
출력의 $V_n$-축을 직접 데이터에 맞추지 않는다($R_n$ 이 축을 이동시키기 때문 — 식별성 \S\ref{sec:ch6_identif}).

\begin{practicebox}
\textbf{inner-loop 안정화(본문 물리식 아님, GS-4$''$).} 초기값 탐색 단계에서 $w_j>0$, $Q_{j,\tot}>0$,
$0\le\chi_j\le1$ 같은 물리 제약을 만족시키려 $\log$/$\mathrm{logit}$ 같은 bounded reparametrization 을 쓸 수
있다. 이는 \emph{최적화 변수의 변환일 뿐} 모델 물리식이 아니며, 최종 보고에서는 원 파라미터값으로 되돌려
신뢰구간과 함께 제시한다. rate 폭주를 임시 clip 하는 것도 디버깅 한정이며 최종 해석에서 제거한다(물리 병목
$k_{j,\lim}$ 으로 대체).
\end{practicebox}

% ====================================================================
\section{피팅 inner-loop 의 수치 엔진 (root-constrained ODE / index-1 DAE)}\label{sec:ch6_dae}
% ====================================================================
\textbf{물리.} 피팅이 한 파라미터 집합에서 모델 ICA 곡선 전체를 얻으려면 $\xi_j(t)$ 의 시간 전개를 풀어야 한다.
정전류 음극 방전에서 $\dd q/\dd t=|I|/Q_\cell$ 이고, 상태변수는 $\xi_j$, $V_n$ 은 전하 보존식이 매 시점 정하는
\emph{algebraic variable} 다. 이는 미분식 1조 $+$ 대수 제약 1조의 \textbf{index-1 DAE}(\AL{60},
\cite{brenan1996,hindmarsh2005})다:
\begin{align}
\frac{\dd\xi_j}{\dd t}&=k_j^{\eff}(V_n,q,|I|,T)\,[\xi_{\eq,j}(V_n)-\xi_j],\label{eq:ch6_ode}\\
0&=F_\cb^\dyn(V_n;q,\bm\xi).\label{eq:ch6_alg}
\end{align}
\textbf{왜 index-1 인가(무비약).} DAE 의 index 는 대수 제약을 미분해 ODE 로 환원하는 데 필요한 미분 횟수다.
식~\eqref{eq:ch6_alg} 을 시간으로 한 번 미분하면 $\dfrac{\partial Q_\bg}{\partial V_n}\dot V_n+\sum_j Q_{j,\tot}
\dot\xi_j-Q_\cell\dot q=0$ 이고, $\partial Q_\bg/\partial V_n\ge\epsilon_Q>0$(식~\eqref{eq:ch6_dyn_deriv})이라
$\dot V_n$ 에 대해 \emph{한 번}에 풀린다 — 따라서 index 가 정확히 1 이다(\AL{60}). 단조성 floor 가 이 풀림의
열쇠이므로, 피팅 중 $\partial Q_\bg/\partial V_n$ 이 0 에 접근하면 DAE 가 ill-conditioned 가 되고 이는
\emph{파라미터가 비물리}라는 신호다(reject; \S\ref{sec:ch6_S1}).

\textbf{시간 적분 절차(정전류; $q(t)$ 명시이므로 root-constrained ODE 로 구현).}
\begin{enumerate}[label=T\arabic*),leftmargin=2.4em]
\item 현재 $q(t),\xi_j(t)$ 에서 식~\eqref{eq:ch6_dyn_root} 로 $V_n(t)$ root-find(S1).
\item $V_n(t)$ 에서 $\xi_{\eq,j}$(S2), $k_j^{\mathrm{act}},k_{j,\lim},k_j^{\eff}$ 또는 $r_j^\pm$(S3) 계산.
\item 식~\eqref{eq:ch6_ode}(또는 forward/backward 형)로 $\dd\xi_j/\dd t$ 계산.
\item implicit solver(BDF 또는 Radau)로 다음 시간으로 진행(\AL{60}, \cite{hindmarsh2005}).
\item 휴지 구간에서는 $q$ 고정, $\xi_j$ relaxation 과 전하 보존 root 를 함께 갱신(Ch2/Ch5 rest; 내부 $I_j\ne0$).
\end{enumerate}
\begin{boundbox}
\textbf{stiff system 처리(cap 아님, GS-4).} 강구동에서 $k_j^{\mathrm{act}}$ 또는 $r_j^+$ 가 매우 커질 수 있다.
이를 arbitrary cap 으로 자르지 않고 \emph{stiff system} 으로 취급한다(implicit BDF/Radau 가 큰 $k_j$ 를 안정적으로
처리; \cite{brenan1996,hindmarsh2005}). 실제 rate 제한은 $k_{j,\lim}$, transport, 유한 전류, site availability
같은 물리적 병목으로 명시한다(Ch1 eq:kphys / Ch3 정합). 즉 \textbf{수치 stiffness 의 해법은 solver 선택이지
모델 변형이 아니다}.
\end{boundbox}
\begin{verifybox}
\textbf{solver tolerance 의 근거(\AL{64}, GS-4$'$).} BDF/Radau 의 상대$\cdot$절대 tolerance $(\mathrm{rtol},
\mathrm{atol})$ 는 임의 상수가 아니라 상태변수의 물리 scale 로 정한다: $\xi_j\in[0,1]$ 이므로 $\mathrm{atol}$ 은
ICA 의 관심 분해능(예: peak 높이 대비 꼬리 기여의 유의 자릿수)으로, $\mathrm{rtol}$ 은 측정 잡음 대비 over-resolve
를 피하는 수준으로 둔다. solver tolerance 를 데이터 잡음보다 훨씬 작게 두는 것은 계산 낭비이고, 훨씬 크게 두면
꼬리(tail-dominated, 작은 기여)를 놓친다 — \textbf{tolerance 는 측정 잡음 floor 와 ICA 분해능 사이}에 둔다.
\end{verifybox}

% ====================================================================
\section{closure: self-consistent Volterra 의 직접 수치 풀이 (기준 해법)}\label{sec:ch6_grid}
% ====================================================================
\textbf{물리.} Ch.1 의 spectrum(eq:spectrum)은 배리어 분포 $\rho_G$ 가 만드는 \emph{연속} relaxation-length 분포다.
실제 풀이는 이를 유한 격자로 이산화한다. 각 배리어 격자점 $g_m$ 의 집단이 독립 1차 완화하고, 전하 보존이 매
시점 $V_n$ 으로 결합한다:
\begin{equation}
\frac{\dd\xi_j(g_m,t)}{\dd t}=k_j^{\eff}(g_m,V_n,q,|I|)\,[\xi_{\eq,j}(V_n)-\xi_j(g_m,t)],\qquad
\xi_j(t)=\sum_{m=1}^{N_g}w_m^{(g)}\,\rho_G(g_m)\,\xi_j(g_m,t),
\label{eq:ch6_grid}
\end{equation}
$\sum_m w_m^{(g)}\rho_G(g_m)\to1$ (격자 정규화; quadrature 가중 $w_m^{(g)}$). \textbf{이 직접 $g$-grid 적분이 본
장의 기준 해법}이다 — 주어진 $\rho_G$ 와 완화 closure 하에서 \emph{이산화 오차만 남는 numerically exact reference}
이며, 모든 축약 해법은 이 해에 대해 검증한다. 이는 Ch.1 self-consistent Volterra(L4)의 직접 수치 평가다.
\begin{verifybox}
\textbf{이중 계상 점검(Ch1 boundbox 의 수치 재확인).} 상수 목표 극한($\Theta_\eq=\Theta_0$, 단일 mode
$A_L=\delta(L-L_0)$, $\Theta_\mathrm{init}=0$)에서 식~\eqref{eq:ch6_grid} 의 이산 해는 $\Theta(q)=\Theta_0
(1-e^{-q/L_0})$ 로 Ch.1 single-mode ODE(eq:single\_kernel)와 일치하고 $q\to\infty$ 서 $\Theta_0$ 로 수렴해야
한다. 만약 kernel 을 \emph{목표} 대신 \emph{잔차}에 곱하면 정상상태가 $\Theta_0/2$ 로 어긋난다 — 수치 구현이
Ch.1 Volterra(목표에 곱)와 같은지 이 극한으로 검증한다.
\end{verifybox}
\begin{boundbox}
\textbf{격자 분해능과 stretched tail(\AL{62}).} Ch.1 의 꼬리는 \emph{tail-dominated}(긴 $L$ mode 가 지배)다.
$g$-격자가 큰 $L$(작은 $k_j$, 큰 배리어) 영역을 충분히 덮지 않으면 꼬리를 놓친다. 따라서 격자 상한은 관심
온도/전위에서 $L_{\max}$ 가 관측 윈도를 덮도록 잡고, 격자 수 $N_g$ 는 $\Theta(q)$ 가 $N_g$ 에 대해 수렴할 때까지
늘린다(자기 수렴 점검; 임의 고정값 아님).
\end{boundbox}
\begin{flagbox}
\textbf{독립 완화 closure = Ch1 이 본 장에 위임한 FLAGGED 항목(미이행).} 식~\eqref{eq:ch6_grid} 의 ``각 $g_m$ 집단이
\emph{독립} 1차 완화''는 Ch.1 eq:xi\_dist 의 평균장(mean-field) ansatz 다. Ch.1 은 이를 flagbox 로 정직 표기하며
``흑연 staging 은 nucleation$\cdot$상경계 이동을 동반하는 \emph{협동적} 1차 상전이라 mode 간 결합(Avrami)을 무시한
평균장이며, 그 타당성은 FLAGGED — Ch.3$\cdot$본 부록(§\ref{sec:ch6_reduction})에서 검증''이라 위임했다. 그러나 위임된 \emph{협동적
staging 검증 자체는 본 작업에서 미이행}이다(독립 완화의 정량 타당성은 operando/calorimetry 측정 + 협동 모형 비교를
요하며 본 이론 범위 밖). 따라서 본 절 $g$-grid 의 ``numerically exact reference''는 그 closure \emph{하에서만}
exact 이며 \emph{exactness 는 이산화 오차에 한하고 closure 자체에는 미치지 않는다} — FLAGGED 로 승계한다.
\end{flagbox}

% ====================================================================
\section{closure 축약: 허용되는 수학적 가속과 그 검증}\label{sec:ch6_reduction}
% ====================================================================
직접 격자(\S\ref{sec:ch6_grid})는 비용이 크므로 피팅 반복에서 가속이 필요할 수 있으나, 물리 의미를 훼손하지
않고 \emph{기준 해 대비 검증}돼야 한다. 다음 축약은 수학적 근거가 있을 때 허용한다.
\begin{enumerate}[label=\textbf{(\alph*)},leftmargin=2.4em]
\item \textbf{Adaptive quadrature.} $\rho_G(g)$ 가 매끄럽고 support 제한($L\ge L_{\min}$)/tail 빠른 감소면, 모델을
  바꾸지 않고 적분 정확도만 확보하는 수치 적분법이므로 허용(\AL{65}).
\item \textbf{Moment matching.} 분포 폭이 작고 kernel 이 매끄러우면 $k_j(g)$/memory kernel 을 저차 moment 로 근사.
  단 moment closure 가 \emph{tail-dominated} kinetics 를 놓치지 않는지 기준 해와 비교한다(Ch.1 stretched tail 주의).
\item \textbf{Prony/rational approximation.} 연속 relaxation spectrum 을 유한 exponential mode 로 근사(\AL{65}):
\begin{equation}
K_j(t)=\int_0^\infty\rho_G(g)\,k_j(g)\,e^{-k_j(g)t}\,\dd g\;\approx\;\sum_{\ell=1}^{N_r}a_\ell\,e^{-t/\tau_\ell},
\qquad a_\ell\ge0,\ \tau_\ell>0.
\label{eq:ch6_prony}
\end{equation}
\end{enumerate}
\textbf{$a_\ell\ge0,\tau_\ell>0$ 의 근거(무비약).} $K_j(t)$ 는 \emph{양의} 분포 $\rho_G k_j$ 위의 지수 감쇠 평균이라
완전 단조(completely monotone) 함수다; 그 유한 mode 근사가 물리적 완화이려면 각 mode 가 \emph{음이 아닌} 진폭과
\emph{양의} 시간상수를 가져야 한다(음수 가중/허수 time constant 는 진동/발산을 낳아 비물리). 따라서 부호 제약은
피팅 편의가 아니라 \emph{완전 단조성의 보존}이다(\AL{65}; KWW stretched 함수와 완화시간 분포의 등가성
\cite{lindsey1980}, 지수 감쇠의 연속합 \cite{johnston2006}). 정전류에서 시간$\leftrightarrow$좌표 대응은
$t-t'\leftrightarrow(q-q')Q_\cell/|I|$, $\tau_\ell\leftrightarrow L_\ell$ 이므로 식~\eqref{eq:ch6_prony} 는 Ch.1
relaxation-length spectrum 의 유한 mode 표현이다.
\begin{boundbox}
\textbf{축약의 한계(되먹임).} $V_n$ 이 전하 보존으로 $\xi_j$ 와 결합돼 고율/큰 lag 에서 $\xi_j\to V_n\to\xi_{\eq,j}
\to\dd\xi_j/\dd t$ 되먹임이 강해진다. Prony/rational/moment 축약은 \emph{준정상$\cdot$약결합} 근사로만 쓰고, 직접
$g$-grid DAE 대비 오차를 \S\ref{sec:ch6_closure} 의 절차로 측정한다. 되먹임 오차가 허용 범위를 넘으면 기준
해법으로 회귀한다.
\end{boundbox}

% ====================================================================
\section{Plan A (Fredholm ratio-substitution closed-form) 의 위치와 ★ 검증 tolerance 의 정직한 정의}\label{sec:ch6_closure}
% ====================================================================
Ch.1 의 \textbf{Plan A}(사용자 PhD ratio-substitution closed-form, \cite{lee2011,son2013,jcp2017})는 문헌적$\cdot$
수학적 기반이 있는 \emph{가속 후보}다(\AL{61}). 본 장은 이를 유지하되, \textbf{직접 $g$-grid 기준 해를 대체하는
기본 모델이 아니라 검증 대상}으로 둔다. Ch.1 의 승격 조건 M1--M5 를 본 장에서 수치적으로 집행한다.
\begin{enumerate}[label=M\arabic*),leftmargin=2.4em]
\item \textbf{kernel mapping.} 흑연 배리어 분포 문제로의 kernel mapping 이 명확하고 적분 안팎 function class 가 같다.
\item \textbf{function class.} ratio-substitution 이 가정하는 분포 class 가 흑연 $\rho_G$ 와 양립한다.
\item \textbf{수렴성.} ratio truncation/Neumann series 의 수렴이 수치적으로 확인된다.
\item \textbf{$\delta$ baseline.} $\rho_G(g)\to\delta(g-g_0)$ 극한에서 단일 배리어 해(Ch.1 single-mode
  $r\propto e^{-(q-q_a)/L_0}$)로 환원된다.
\item \textbf{정량 오차.} 기준 해 대비 상대오차가 검증 tolerance 이하다(아래 ★ 정의).
\end{enumerate}

\subsection{★ 검증 tolerance $\varepsilon_\tol$ 의 정직한 재정의 (방향-3: FLAGGED 허위 정정)}\label{sec:ch6_epsilon}
\textbf{문제(Phase A 적발).} consolidated\_ch6 의 M5 는 정량 오차 조건을
\[
\varepsilon=\frac{\lVert\Theta_{\mathrm A}-\Theta_{\mathrm B}\rVert}{\lVert\Theta_{\mathrm B}\rVert}\le\varepsilon_\tol
\]
로 쓰면서 $\varepsilon_\tol$ 을 \emph{established 임계값처럼} 사용했다. 그러나 (i) 그 \emph{값}의 근거가 본문에
없고, (ii) 문서 전체에 정직 표기용 flagbox 가 0건이었다. 이는 GS-3$\cdot$GS-4 위반(근거 없는 수치 임계값)이다.
본 절은 이를 \textbf{근거 있는 수치 tolerance 로 재서술}하되, 정당화되지 않는 특정 숫자는 FLAGGED 로 정직 표기한다.

\textbf{무엇을 비교하는가(정의 명확화).} $\Theta_{\mathrm B}$ 는 직접 $g$-grid 기준 해(\S\ref{sec:ch6_grid}),
$\Theta_{\mathrm A}$ 는 Plan A closed-form 해다. norm $\lVert\cdot\rVert$ 은 \emph{관심 tail-window} 위의 이산
$L^2$ 노름으로 명시한다: $\lVert f\rVert^2=\sum_{i\in\mathcal W}(f(q_i))^2\,\Delta q$, $\mathcal W=$ post-peak 꼬리
구간. window 를 명시하지 않으면 peak 영역의 큰 값이 분모를 지배해 꼬리 오차를 가린다(\textbf{꼬리가 본 이론의
관심사}이므로 window 가 필수다).

\textbf{tolerance 의 \emph{근거 있는} 결정(임의 상수 금지).} $\varepsilon_\tol$ 은 ``작으면 좋다''가 아니라 두 물리
scale 의 \emph{큰 쪽}으로 둔다(\AL{66}):
\begin{equation}
\varepsilon_\tol\;=\;\max\!\Big(\underbrace{\varepsilon_{\mathrm{disc}}}_{\text{기준 해 자체의 이산화 오차}},\
\underbrace{\varepsilon_\noise}_{\text{측정 잡음으로 인한 ICA 불확실도}}\Big).
\label{eq:ch6_epsilon}
\end{equation}
\emph{한 줄씩 근거}: (a) Plan A 를 기준 해 $\Theta_{\mathrm B}$ 보다 더 정확하게 만들 수는 없다 — $\Theta_{\mathrm B}$
자신이 격자 이산화 오차 $\varepsilon_{\mathrm{disc}}$ 를 품기 때문이다. 따라서 $\varepsilon_\tol<\varepsilon_{\mathrm{disc}}$
를 요구하는 것은 무의미하다(기준 해의 정밀도를 넘는 일치를 강요). $\varepsilon_{\mathrm{disc}}$ 는 \S\ref{sec:ch6_grid}
의 격자 자기수렴(격자 배가 시 $\Theta_{\mathrm B}$ 변화)으로 \emph{측정}한다. (b) 모델 예측을 데이터 잡음보다 더
정밀하게 일치시키는 것은 관측 가능한 차이가 아니다 — ICA 의 잡음 유발 불확실도 $\varepsilon_\noise$ 는 측정
전압/전류 잡음을 $\dd Q/\dd V$ 로 전파해 추정한다($dQ/dV$ 는 미분이라 잡음이 증폭됨, \cite{dubarry2022,fly2020}
\AL{63}). 따라서 \textbf{Plan A 가 기준 해와 ``측정으로 구별 불가능한 만큼'' 일치하면 충분}하며, 그 경계가
식~\eqref{eq:ch6_epsilon} 다. 두 scale 중 \emph{큰 쪽}을 쓰는 이유는, 더 작은 쪽까지 맞추려는 것이 위 (a)(b)
어느 한쪽에서 무의미해지기 때문이다.
\begin{flagbox}
\textbf{$\varepsilon_\tol$ 의 \emph{특정 숫자값}은 미확정(\AL{66}).} 식~\eqref{eq:ch6_epsilon} 은 $\varepsilon_\tol$
을 \emph{측정 가능한 두 scale} 로 환원했으나, $\varepsilon_{\mathrm{disc}},\varepsilon_\noise$ 의 \emph{수치}는
사용자 보유 데이터의 실제 잡음 수준$\cdot$격자 수렴 거동에서 산출돼야 한다 — 본 이론 문서가 ``$\varepsilon_\tol=10^{-2}$''
같은 \emph{고정 숫자}를 established 로 제시하지 않는다. consolidated\_ch6 가 $\varepsilon_\tol$ 을 무근거 임계값처럼
쓴 것을 본 절이 \emph{측정 기반 정의}로 교체했고, 구체 값은 데이터 확보 후 결정하는 \textbf{미확정 항}으로 정직
표기한다(GS-3$\cdot$GS-4$'$ 준수). $L^2$ vs $L^\infty$ norm 선택, window 경계 정의도 데이터 의존이라 함께 미확정이다.
\end{flagbox}
\begin{enumerate}[label=F\arabic*),leftmargin=2.4em,start=6]
\item 실패 시 fallback 순서: 직접 grid(\S\ref{sec:ch6_grid}) $\to$ adaptive quadrature $\to$ Prony/rational
  (\S\ref{sec:ch6_reduction}). 즉 \emph{기준 해는 언제나 직접 grid} 이며 Plan A 는 통과 시에만 가속용으로 끼운다.
\end{enumerate}
\begin{boundbox}
\textbf{Plan A 의 broad-kernel 한계(★ \AL{62}, jcp2017 자기명시).} 사용자 논문(\cite{jcp2017}, Eq.~34 직후)은
``ratio-substitution 근사의 정확도는 reaction zone(분포)이 매우 넓어질 때 나빠진다''고 \emph{스스로 경고}한다.
graphite 의 넓은 spectrum($=$저온 stretched tail, \emph{관심 영역})은 바로 그 broad-kernel 영역이라 Plan A 가
\textbf{가장 부정확할 수 있다}. 또한 사용자 논문은 \emph{Fredholm}(고정 구간) ratio closure 이고 graphite 문제는
\emph{Volterra}(인과 상한 $q$, Ch.1 L4)라 자기참조 선형화가 선행돼야 한다. 따라서 \textbf{Plan A 단독 채택 금지;
항상 직접 grid 대비 식~\eqref{eq:ch6_epsilon} 검증.} graphite transition 개수$\cdot$분포 폭$\cdot$hysteresis branch
offset 을 Fredholm correction factor 안에 흡수하지 않는다(식별성 보존).
\end{boundbox}

% ====================================================================
\section{파라미터 식별성(identifiability)과 공선형 차단}\label{sec:ch6_identif}
% ====================================================================
\textbf{물리.} 피팅의 핵심 위험은 \emph{여러 파라미터가 같은 곡선 변화를 만들어 서로 흉내 내는} 공선형
(co-linearity)이다. 흑연 ICA 에서 가장 위험한 두 쌍을 먼저 닫고, 일반 규칙을 제시한다.

\subsection{$R_n$ vs $\chi_j$ — 분극과 mobility 가속의 공선형}\label{sec:ch6_RnChi}
\textbf{둘 다 꼬리를 늘린다.} 분극 $R_n$ 은 apparent 전위축을 이동시켜(eq:Vapp) ICA peak 를 넓히고, mobility 가속
$\chi_j$ 는 $L=|I|/(Q_\cell k_j)$ 를 통해 꼬리 길이를 직접 바꾼다(Ch.1 ideal-width 형
$\partial\ln L/\partial V_{n,\drive}=-\chi_j s_{\phi,j}F/RT$; 일반형은 $-\chi_j s_{\phi,j}n_j^{\eff}F/RT
=-\chi_j s_{\phi,j}/w_j$, $n_j^{\eff}=RT/Fw_j$). 한 데이터에서 둘을 동시에 자유 피팅하면 공선형이라 분리 불가다(Ch1 \S falsify, Ch3
\AL{39}).

\textbf{분리 전략(순차 제약).} \emph{독립 실험으로 $R_n$ 을 먼저 고정}한다: 펄스/전류 차단(GITT 의 순간 step)은
$|I|\to0$ 직후 전압 도약으로 \emph{순간} 분극을 주며, 이는 mobility 완화(느린 spectrum)와 시간척도가 분리된다.
$R_n$ 을 이렇게 제한한 \emph{뒤에} 같은 모델로 $\chi_j$ 를 본다. 순서를 바꾸면(먼저 $\chi_j$ 자유화) 분극이
$\chi_j$ 로 흡수돼 가짜 가속을 얻는다.
\begin{keybox}
\textbf{세 저항은 다른 물리다(Ch1$\cdot$Ch3$\cdot$Ch4 계승).} apparent voltage-axis 계수 $R_n$ $\ne$ charge-transfer
저항 $R_\ct=RT/(FA_\eff j_{0,n})$(Ch3 small-signal) $\ne$ heat-generation 저항 $R_{n,\eff}$(Ch4 발열). 피팅에서
이 셋을 한 ``$R$''로 합치면 식별성과 발열 예측이 모두 무너진다 — \textbf{$R_n\ne R_\ct\ne R_{n,\eff}$} 를
파라미터 단계에서 분리 유지한다.
\end{keybox}

\subsection{$R_n/R_\ct/R_{n,\eff}$ 의 실험적 분리 (Ch3 계승)}\label{sec:ch6_threeR}
$R_\ct$ 는 \emph{소진폭 EIS/Tafel}(small-signal 선형 영역, \cite{bardfaulkner})에서, $R_n$ 은 \emph{ICA 전압축
이동}(apparent), $R_{n,\eff}$ 는 \emph{calorimetry}(발열률 $\div I^2$)에서 각각 독립 제약한다. 한 실험이 셋을 모두
주지 않으므로, 같은 데이터에서 동시에 자유화하지 않는다(\AL{39},\AL{47},\AL{68}).

\subsection{일반 식별성 규칙 — 동시 자유화 금지 집합}\label{sec:ch6_noFree}
다음을 \emph{같은} 자료에서 동시에 자유화하면 식별성이 무너진다(Ch1--3 누적 경고의 통합):
\begin{equation}
\{\,R_n,\ k_j^{\eff},\ w_j,\ \rho_G(g),\ Q_\bg,\ Q_{j,\tot},\ k_{j,\lim}\,\}.
\label{eq:ch6_no_free}
\end{equation}
\textbf{순차 제약 사슬(deliverable).} 식~\eqref{eq:ch6_no_free} 의 공선형은 \emph{독립 실험을 순서대로 거는} 것으로만
풀린다. 본 장의 권고 순서는 \S\ref{sec:ch6_ocvfit}--\ref{sec:ch6_crate} 에서 단계별로 실현한다:
\[
\underbrace{\text{저속 OCV}}_{U_j,w_j,Q_{j,\tot},Q_\bg}\ \to\ \underbrace{\text{GITT}}_{R_n,\ \text{완화 signature}}\ \to\
\underbrace{\text{다온도 Arrhenius}}_{\Delta H_a,\Delta S_a,k_0}\ \to\ \underbrace{\text{C-rate series}}_{\chi_j,\ k_{j,\lim}}.
\]
$\rho_G$ 는 어느 단계에서도 관측 꼬리에서 \emph{역산하지 않고}(비유일, \AL{16}) 독립 입력으로만 둔다.

% ====================================================================
\section{S1 실현: 저속 OCV 로 열역학 골격 제약}\label{sec:ch6_ocvfit}
% ====================================================================
저속(준평형) OCV/STO--OCV 테이블을 평형 음극 전위 $V_{n,\OCV}(q)$ 의 실측 proxy 로 쓴다. \textbf{이 자료만으로
$U_j,w_j,Q_{j,\tot},Q_\bg$ 가 유일 결정되지는 않으므로} 다음 물리 제약 하에서 피팅한다(과적합 차단):
\begin{enumerate}[label=O\arabic*),leftmargin=2.4em]
\item 총용량 정합: $\sum_j Q_{j,\tot}+\Delta Q_\bg=Q_\cell\cdot\Delta q_{\mathrm{full}}$(Ch1 전하 보존의 적분형).
\item $Q_\bg(V_n)$ 단조성 $\partial Q_\bg/\partial V_n\ge\epsilon_Q>0$ 및 smoothness(Ch1 eq:monotone).
\item transition 개수 $N_p$ 를 필요 이상 늘리지 않음(parsimony; 관측 $N_p\approx3\sim4$, \AL{3}).
\item $Q_{j,\tot}>0$(방전 방향 용량 기여; Ch1).
\item $w_j$ 는 \emph{평형} transition width 이며 kinetic 배리어 분포 $\rho_G(g)$ 와 혼동 금지(Ch1 구분; 단
  $n_j^{\eff}w_j=RT/F$ 정합, Ch3 \AL{34}).
\item 잔차를 임의 background wiggle 로 모두 흡수하지 않음(GS-4; 잔차는 다음 단계가 설명할 신호일 수 있음).
\end{enumerate}
\textbf{평형 형태의 판별(Ch1 N1 의 실현).} 저속 OCV 의 near-peak 감쇠를 $\log(\dd Q/\dd V)$ 대 $(V-U)$(직선이면
logistic, Ch1 eq:logistic) 또는 대 $(V-U)^2$(직선이면 Gaussian-cumulative)로 그려 평형 isotherm 형태를
\emph{판별}한다. \textbf{logistic 이외 형태(erf 등)는 흑연에 대한 열역학적 유도가 없는 FLAGGED ansatz}(Ch1
flagbox)이므로, 데이터가 logistic 을 기각하면 그때 BOUNDED ansatz 로만 도입한다.
\begin{boundbox}
\textbf{저속이 평형은 아니다.} $0.1/0.2\,\mathrm C$ 는 \emph{quasi-equilibrium proxy} 이지 strict thermodynamic
equilibrium 이 아니다(\AL{63}). peak 위치/폭에 잔존 kinetic broadening 이 섞일 수 있으므로, OCV 골격은 GITT
relaxation 끝점(\S\ref{sec:ch6_gittfit})과 교차 검증한다.
\end{boundbox}

% ====================================================================
\section{S1--S2 보강: GITT 로 분극과 완화 signature 제약}\label{sec:ch6_gittfit}
% ====================================================================
GITT(galvanostatic intermittent titration)는 짧은 정전류 펄스 $+$ 긴 휴지의 반복으로 \emph{순간 응답}과 \emph{완화
응답}을 함께 준다. \textbf{GITT relaxation 을 곧바로 $k_j$ 의 직접 측정값으로 해석하지 않는다} — 여러 과정이
섞이기 때문이다:
\begin{equation}
\text{GITT relaxation}=\underbrace{\text{solid diffusion}}_{}+\underbrace{\text{electrolyte relaxation}}_{}+
\underbrace{\text{phase/staging relaxation}}_{\text{Ch1 }k_j}+\underbrace{\text{OCV curvature}}_{}+\underbrace{\text{measurement filtering}}_{}.
\label{eq:ch6_gitt_decomp}
\end{equation}
\textbf{권장 절차(순차 분리).}
\begin{enumerate}[label=G\arabic*),leftmargin=2.4em]
\item 펄스 직후 \emph{순간} step 으로 ohmic$\cdot$fast polarization 제약($R_n$ 의 일부; Ch3 $\eta_\ct$ 분리 시작).
\item 장시간 완화 끝점을 저속 OCV(\S\ref{sec:ch6_ocvfit})와 비교해 평형 골격 검증(O 제약과 정합 확인).
\item relaxation curve 를 단일 $k_j$ 로 단정하지 않고 single- vs multi-time-constant 후보를 비교한다(Ch.1
  spectrum 의 다중 $L$ — 여러 시간상수가 보이면 분포 $\rho_G$ 의 직접 증거).
\item C-rate series(\S\ref{sec:ch6_crate})와 결합해 $k_j$, diffusion limitation, transport limitation 을 분리.
\end{enumerate}
\textbf{$\chi_j$ vs $\eta_\ct$ 공선형(Ch1/Ch3 N 의 실현).} GITT 의 순간 step($\eta_\ct$, charge-transfer)과 느린
relaxation(barrier spectrum)을 분리한 \emph{후에만} $\chi_j(=\beta_j$; 조건부, 대칭 Marcus$\cdot$정상 target 극한 한정$)$ 귀속이 성립한다. 즉 \textbf{$R_n$ 과 $\chi_j$
는 GITT 의 시간척도 분리로 비로소 독립 제약}된다(\S\ref{sec:ch6_RnChi} 의 분리 전략을 데이터로 실현).

% ====================================================================
\section{S3 핵심: 다온도 Arrhenius 로 활성화 파라미터 제약}\label{sec:ch6_arrhenius}
% ====================================================================
\textbf{물리.} 활성화 파라미터 $\Delta H_{a,j},\Delta S_{a,j},k_0$ 는 \emph{단일 온도} 데이터로는 분리되지 않는다 —
$k_j=k_0 e^{-(\Delta H_a-T\Delta S_a-\chi_j\mathcal A_j)/RT}$ 에서 한 온도는 한 방정식만 주기 때문이다. 여러 온도가
있어야 Arrhenius 직선의 \emph{기울기와 절편}이 따로 정해진다.

\textbf{Arrhenius 분해(무비약).} 꼬리 길이 $L=|I|/(Q_\cell k_j)$ 의 로그를 취하면(Ch1 eq:L\_of\_G)
\begin{equation}
\ln\frac1L\;=\;\ln k_j-\ln\frac{|I|}{Q_\cell}\;=\;\ln k_0-\frac{\Delta H_{a,j}}{RT}+\frac{\Delta S_{a,j}}{R}
+\frac{\chi_j\mathcal A_j}{RT}-\ln\frac{|I|}{Q_\cell}.
\label{eq:ch6_arrhenius}
\end{equation}
\emph{한 줄씩}: $\Delta G_{a,j}=\Delta H_{a,j}-T\Delta S_{a,j}$(Ch1 \AL{10})를 $k_j=k_0 e^{-(\Delta G_{a,j}-\chi_j
\mathcal A_j)/RT}$ 에 넣고 로그를 취했다. \textbf{$\mathcal A_j$ 를 상수로 고정하고}(같은 동작점) $\ln(1/L)$ 을 $1/T$ 에
대해 그리면, $\chi_j\mathcal A_j/(RT)=(\chi_j\mathcal A_j/R)\cdot(1/T)$ \emph{도} $1/T$ 에 선형이라 기울기에 들어간다.
따라서 기울기는 순수 $-\Delta H_{a,j}/R$ 가 \emph{아니라}
\begin{equation}
\frac{\dd\ln(1/L)}{\dd(1/T)}\Big|_{\mathcal A_j}=-\frac{\Delta H_{a,j}-\chi_j\mathcal A_j}{R}=-\frac{\Delta H_{a,j}^{\eff}}{R}
\label{eq:ch6_arrhenius_slope}
\end{equation}
즉 \emph{유효} 활성화 엔탈피 $\Delta H_{a,j}^{\eff}=\Delta H_{a,j}-\chi_j\mathcal A_j$ 를 준다. \textbf{순수 $\Delta H_a$ 를
얻으려면} 구동력을 없앤 극한 $\mathcal A_j\to0$(꼬리 영역 $V_{n,\drive}\to U_j$, 무구동)로 외삽하거나, $\chi_j$ 가 독립
제약된 뒤 $\chi_j\mathcal A_j/R$ 만큼 보정한다. $1/T\to0$ 외삽 절편 $=\ln k_0+\Delta S_{a,j}/R-\ln(|I|/Q_\cell)$
(prefactor$\cdot$엔트로피 묶음; $\ln(|I|/Q_\cell)$ 상수항 포함)다. $k_0$ 와 $\Delta S_a$ 는 절편에서 공선형이므로
$k_0=(k_BT/h)\kappa$ 의 transition-state 절대값을 \emph{인용}(Ch1 \AL{10}, \cite{eyring1935})해 $\Delta S_a$ 를
분리하되, $\kappa(T)$ 미지로 $\ln\kappa$ 가 $\Delta S_a$ 로 흡수되는 잔여 공선형은 BOUNDED 다(\AL{69};
\cite{bardfaulkner}). $\Delta S_a$ 를 reaction entropy(Ch2 \AL{24}, 별개 물리)와 혼동하지 않는다.
\begin{boundbox}
\textbf{$\mathcal A_j$ 고정의 의미(★ 정정).} 흔한 오해와 달리, $\chi_j\mathcal A_j/(RT)$ 항은 $\mathcal A_j$ 를
\emph{고정해도} $1/(RT)=(1/R)(1/T)$ 를 통해 기울기에 기여한다 — 따라서 고정 $\mathcal A_j$ Arrhenius 기울기는 순수
$-\Delta H_a/R$ 가 아니라 \emph{유효} $-(\Delta H_a-\chi_j\mathcal A_j)/R$ 다(식~\eqref{eq:ch6_arrhenius_slope}).
전위를 \emph{바꾸면} $\mathcal A_j$ 자체가 변해 기울기를 더 휘게 하므로, 깨끗한 분리는 한 동작점 기울기에서
$\chi_j\mathcal A_j/R$ 를 보정하거나 $\mathcal A_j\to0$ 외삽으로 한다. 이 $\chi_j\mathcal A_j$ 혼입이 N1
반증(\S\ref{sec:ch6_falsify})의 대상이다. (참고: $\mathcal A_j$ 를 일반형 $n_j^\eff F(V-U)=(RT/w_j)(V-U)$ 로 쓰면
$\chi_j\mathcal A_j/RT=\chi_j(V-U)/w_j$ 로 explicit $1/T$ 가 상쇄되나, 그땐 $w_j(T)$ 의 온도 표류가 같은 자리에 남아
— 어느 convention 이든 순수 $\Delta H_a$ 분리는 $\mathcal A_j\to0$/보정이 필요하다.)
\end{boundbox}
\textbf{Arrhenius 직선성 = 반증 가능성(Ch1 N2).} $\ln(1/L)$ 대 $1/T$ 가 \emph{직선이 아니면} 단일 활성화 그림이
기각된다 — 분포 $\rho_G$ 의 폭이 온도에 따라 유효 기울기를 휘게 하거나(stretched, Ch1 \AL{14}), transport 가
지배하는 경우다. 직선성 자체가 모델의 검정이다.

% ====================================================================
\section{S3--S5 검증: C-rate series 로 전위 의존 가속과 꼬리 제약}\label{sec:ch6_crate}
% ====================================================================
$0.1/0.2\,\mathrm C$ 는 quasi-equilibrium proxy(peak 위치/폭), $1.0\,\mathrm C$ 는 kinetic-lag/transport stress
test 다(strict frozen limit 아님). \textbf{검증 순서(순차 제약의 마지막 단계).}
\begin{enumerate}[label=C\arabic*),leftmargin=2.4em]
\item 저속 OCV 로 평형 골격($U_j,w_j,Q_{j,\tot},Q_\bg$) 제약(\S\ref{sec:ch6_ocvfit}).
\item GITT 로 fast polarization($R_n$)$\cdot$relaxation signature 제약(\S\ref{sec:ch6_gittfit}).
\item 다온도로 활성화 파라미터($\Delta H_a,\Delta S_a,k_0$) 제약(\S\ref{sec:ch6_arrhenius}).
\item $0.1/0.2\,\mathrm C$ 에서 ICA/DVA peak 위치$\cdot$폭 검증(S1--S2; 평형 골격 재확인).
\item $0.5/1.0\,\mathrm C$ 에서 \textbf{꼬리$\cdot$broadening$\cdot$shift} 로 $\chi_j$ 제약(S3--S6, kernel integral).
  전위 의존 $\partial\ln L/\partial V_{n,\drive}<0$ 가 C-rate 증가 시 꼬리 단축으로 나타나는지 확인(\AL{62}).
\item 잔차를 $\rho_G(g),k_{j,\lim},R_n$, transport 중 \emph{어느 항으로} 설명 가능한지 독립 제약과 함께 판단
  (한 항으로 모두 흡수 금지; 식~\eqref{eq:ch6_no_free}).
\end{enumerate}
\textbf{이 순서가 식별성을 보장하는 이유.} 각 단계가 \emph{앞 단계에서 고정된 파라미터} 위에서 \emph{새 파라미터
하나(군)} 만 자유화하므로, 식~\eqref{eq:ch6_no_free} 의 공선형 집합이 한 번에 풀리지 않고 차례로 닫힌다 — 이것이
``동시 자유 피팅'' 과 ``순차 식별'' 의 결정적 차이다.

% ====================================================================
\section{반증(falsification): forward-only 예측의 데이터 검정}\label{sec:ch6_falsify}
% ====================================================================
\textbf{물리.} Ch.1--5 는 모두 \emph{forward-only}(파라미터$\to$예측) 모델이다. 피팅이 ``곡선을 맞췄다''로 끝나면
어떤 모델도 충분한 자유도로 맞출 수 있어 의미가 없다. \textbf{모델의 가치는 \emph{반증 가능성}}이며, 본 절은 Ch.1
의 반증 규칙 N1--N4 를 \emph{데이터 시험}으로 실현한다.
\begin{loopbox}
\textbf{도입 관찰의 반증 가능 시험(Ch1 N1--N4 의 데이터 실현).} 도입 관찰 ``저온$\cdot$고율서 긴 겹침 꼬리''를
다음으로 시험한다(어느 하나라도 어긋나면 해당 귀속을 \emph{기각}).
\begin{itemize}[leftmargin=1.6em]
\item \textbf{(N1) Arrhenius 기울기 vs $E_D$.} 꼬리 길이의 $\ln(1/L)$ 대 $1/T$ 기울기(식~\eqref{eq:ch6_arrhenius})가
  독립 측정 확산 활성화 $E_D$ 와 \emph{일치}하고 \emph{전위 의존이 없으면}, 꼬리는 barrier-spectrum 이 아니라
  transport 지배다 $\Rightarrow$ barrier-spectrum 귀속 \emph{기각}. C-rate$\times$온도 교차표로 측정.
\item \textbf{(N2) 전위 의존 spectrum shift.} 전위/$|I|$ 증가 시 spectrum 이 short-$L$ 로 이동하는지
  ($\partial\ln L/\partial V_{n,\drive}<0$, Ch1 eq:L\_of\_G). 이동이 없으면 $\chi_j$ 항 \emph{기각}(가속의 전위
  의존이 모델의 핵심 예측).
\item \textbf{(N3) 저온 lineshape.} 저속 OCV near-peak lineshape 이 저온에서 \emph{안 좁아지면}, 평형 broadening
  ($w=RT/F$ 는 저온서 좁아짐, Ch1 eq:dxidV)으로 설명 안 되는 heterogeneous disorder 기여가 존재 $\Rightarrow$
  kinetic 단독 귀속 기각, disorder 와 분리(Ch1 flagbox).
\item \textbf{(N4) $\rho_G$ forward-only.} $\rho_G(g)$ 를 관측 꼬리에서 \emph{역산하지 않고}(비유일, \AL{16})
  독립 근거(입자 크기 분포$\cdot$EIS$\cdot$계산)로 주어 forward 예측만 한다 — 역산하면 비유일이라 반증력이 없다.
\end{itemize}
이로써 \textbf{도입 관찰 $\to$ 인과 사슬(Ch1) $\to$ 발열/속도론/히스테리시스(Ch2--5) $\to$ 피팅식 $\to$ 반증
가능한 데이터 시험(본 부록)} 으로 한 바퀴가 닫힌다.
\end{loopbox}
\begin{boundbox}
\textbf{$\rho_G$ 역산 금지의 피팅 함의(\AL{16}).} 피팅 도구는 잔차를 줄이려 $\rho_G$ 를 자유화하고 싶어 하지만,
stretched 거동에서 $\rho_G$ 로의 역매핑은 비유일(재포획 등 다른 메커니즘도 같은 꼬리)이다. 따라서 $\rho_G$ 를
피팅 변수로 풀면 반증력이 사라진다 — 독립 입력으로 고정하고, 그 입력이 데이터를 \emph{예측}하는지만 본다.
\end{boundbox}

% ====================================================================
\section{calorimetry 부재 시 발열 항의 forward-prediction 제한}\label{sec:ch6_heat}
% ====================================================================
\textbf{Ch2/Ch4 의 발열 항을 ICA/전압 데이터만으로 피팅하지 않는다}(\AL{67}). 비가역 발열
$\dot{\mathcal Q}_\irr=\sum_j n_j^{\eff}I_j\eta_j\ge0$(Ch4)과 가역 발열 $\dot{\mathcal Q}_\rev$(Ch2 Bernardi
energy balance \cite{bernardi1985}; 다공성 전극 열모델 \cite{thomasnewman2003})은 \emph{전류$\cdot$과전압$\cdot$
entropy coefficient} 로부터 \textbf{forward 로 예측}만 하고, 그 예측을 \emph{calorimetry} 가 확보됐을 때
검증한다(\AL{67}). 이유: ICA/전압 데이터는 발열을 직접 보지 않으므로, 발열 파라미터를
전압 잔차에 피팅하면 다른 전압 효과(분극 등)와 공선형이 되어 가짜 값을 얻는다(식~\eqref{eq:ch6_no_free} 의
정신을 발열로 확장).
\begin{boundbox}
\textbf{발열의 ICA 되먹임(등온 한계의 신호).} 등온 가정이 깨지면(고율 자기발열) Ch.2 의 되먹임
($\dot{\mathcal Q}\to T\uparrow\to k_j\uparrow\to L\downarrow$ 꼬리 단축, 동시에 $w=RT/F\uparrow$ 로 평형 peak
폭 증가)이 ICA tail 에 두 부호로 작용한다. 따라서 고율 ICA tail 을 단일 부호로 단정하지 않고, 등온이 의심되면
Ch2/Ch4 thermal coupling 으로 돌아간다(\AL{67}).
\end{boundbox}

% ====================================================================
\section{충전$\cdot$양극$\cdot$full-cell 로의 확장 조건 (음극 방전 검증 후)}\label{sec:ch6_extension}
% ====================================================================
\begin{enumerate}[label=E\arabic*),leftmargin=2.4em]
\item \textbf{충전 branch.} Ch.5 의 signed current $I_s$ 와 metastable branch target $\xi_{\tar,j}^b=\xi_{\sstat,j}^b$,
  유도된 부호 $s_{\phi,j}^b=-1$ 을 사용한다(\AL{53}--56). branch switching 시 $\xi_j,h_j,V_n$ continuity(재초기화 X),
  branch 별 root finding(\S\ref{sec:ch6_S1})으로 매 시점 $V_n$ 결정.
\item \textbf{히스테리시스 분리 검증.} $\Delta V_\obs\ne\Delta V_\hys$(Ch5 \AL{58}): 관측 loop 폭에서 polarization
  ($R_n^b$; $|I_s|\to0$ 외삽으로 소멸)$\cdot$kinetic lag$\cdot$transport 를 빼고 남는 잔존 분지 차만 true
  hysteresis 후보로 둔다. empirical branch fit 은 metastable free-energy/domain memory 분리 안 되면 \emph{validation
  candidate} 로만 취급(주 결론서 제외, Ch5 forward-only).
\item \textbf{양극.} 양극에 별도 전하 보존식$\cdot$STO--OCV$\cdot$GITT 를 같은 절차로 적용.
\item \textbf{full-cell.} $V_\cell=V_p-V_n-\eta_\loss$ 또는 apparent electrode voltage 표현; loss/polarization
  이중 계상 금지(Ch5). full-cell ICA/DVA 는 양극/음극/branch signed loss/transport/polarization 이 섞인 관측
  신호라 음극 단독과 직접 같지 않다.
\item \textbf{순서.} \textbf{음극 방전이 먼저 식별$\cdot$검증되지 않으면 full-cell 조합 해석으로 넘어가지 않는다.}
\end{enumerate}

% ====================================================================
\section{금지되는 피팅 편의와 실무 보조의 격리}\label{sec:ch6_forbidden}
% ====================================================================
\textbf{금지(본문 물리식으로 쓰지 않음, GS-4).}
\begin{enumerate}[label=(\arabic*),leftmargin=2.4em]
\item 강구동 rate 폭주를 물리 의미 없는 cap/clip 으로 자르기(대신 stiff solver $+$ 물리 병목 $k_{j,\lim}$).
\item 잔차를 arbitrary background correction(임의 wiggle)으로 모두 흡수(O6).
\item branch hysteresis 를 단순 charge/discharge offset 으로만 피팅(Ch5 metastable 분리 없이).
\item negative quadrature weight/비물리 time constant 로 relaxation spectrum 맞추기(완전 단조성 위반, \S\ref{sec:ch6_reduction}).
\item Fredholm/Padé(Plan A) correction factor 를 불분명한 fitting knob 으로 사용(\S\ref{sec:ch6_closure}).
\item $\rho_G$ 를 관측 꼬리에서 역산(비유일, N4).
\item 근거 없는 수치 임계값($\varepsilon_\tol$ 등)을 established 로 사용(\S\ref{sec:ch6_epsilon} 의 측정 기반 정의로 대체).
\end{enumerate}
\begin{practicebox}
\textbf{실무 보조(본문 물리식 아님, GS-4$''$).} 초기값 탐색$\cdot$디버깅에서 bounded transform, 임시 clipped rate,
smoothed background 같은 수치 보조를 \emph{임시} 사용할 수 있으나 본 모델 물리식이 아니다. 최종 해석$\cdot$논문용
식$\cdot$검증 결과에서는 제거하고 물리적 제한항(병목 $k_{j,\lim}$) 또는 검증된 축약법으로 대체한다. 피팅 보조가
결과에 남았는지 여부는 \S\ref{sec:ch6_selfcheck} 검수표가 점검한다.
\end{practicebox}

% ====================================================================
\section{부록 B 자체 검수표 (구 Chapter 6)}\label{sec:ch6_selfcheck}
% ====================================================================
{\small
\begin{longtable}{p{0.74\textwidth} p{0.16\textwidth}}
\toprule 검수 항목 & 결과 \\ \midrule \endhead
Ch1--5 본문 수정 없이 부록만 적층(P3-6, P5) & 통과(\S\ref{sec:ch6_inherit}) \\
★ 피팅 실무가 주역, 수치해법은 inner-loop solver 로 종속(사용자 5-30) & 통과(\S\ref{sec:ch6_S},\ref{sec:ch6_dae}) \\
피팅 평가 순서 S1--S6 을 피팅 워크플로로 실현 & 통과(\S\ref{sec:ch6_S}) \\
피팅 파라미터 전수 목록 $+$ 정의 장$\cdot$1차 제약 실험 매핑 & 통과(\S\ref{sec:ch6_inherit},\ref{sec:ch6_identif}) \\
★ 식별성: $R_n$ vs $\chi_j$ 공선형 분리(독립 실험 순차) & 통과(\S\ref{sec:ch6_RnChi}) \\
★ $R_n\ne R_\ct\ne R_{n,\eff}$ 분리 유지(Ch1$\cdot$3$\cdot$4 계승) & 통과(\S\ref{sec:ch6_threeR} keybox) \\
동시 자유화 금지 집합 $+$ 순차 제약 사슬(OCV$\to$GITT$\to$Arrhenius$\to$C-rate) & 통과(\S\ref{sec:ch6_noFree}) \\
dynamic root vs OCV root 의 derivative 구분(식~\eqref{eq:ch6_dyn_deriv} vs \eqref{eq:ch6_ocv_root}) & 통과(\S\ref{sec:ch6_S1}) \\
OCV root 분모 $=$ Ch2 평형 chemical capacitance 정합 & 통과(\S\ref{sec:ch6_S1}) \\
전하 보존식 해 존재 조건을 residual 로 흡수 X(비물리 파라미터 reject) & 통과(\S\ref{sec:ch6_S1},\ref{sec:ch6_dae}) \\
index-1 DAE \emph{왜 index-1 인지} 유도(단조성 floor) & 통과(\S\ref{sec:ch6_dae}) \\
강구동 rate 를 cap 으로 자르지 않고 stiff$+$물리 병목으로 처리 & 통과(\S\ref{sec:ch6_dae} boundbox) \\
직접 $g$-grid 를 기준 해법으로; 모든 축약은 대비 검증 & 통과(\S\ref{sec:ch6_grid}) \\
이중 계상(목표 vs 잔차) 수치 점검 극한 & 통과(\S\ref{sec:ch6_grid} verifybox) \\
Prony/rational $a_\ell\ge0,\tau_\ell>0$ 을 완전 단조성으로 유도(편의 아님) & 통과(\S\ref{sec:ch6_reduction}) \\
★ FLAGGED 허위 정정: $\varepsilon_\tol$ 을 측정 기반(이산화 오차$\vee$잡음)으로 재정의 & 통과(\S\ref{sec:ch6_epsilon} eq~\eqref{eq:ch6_epsilon}) \\
★ $\varepsilon_\tol$ 특정 숫자값을 flagbox 로 정직 표기(established 사용 X) & 통과(\S\ref{sec:ch6_epsilon} flagbox) \\
★ root-find/solver tolerance 도 물리 scale$\cdot$잡음으로 근거화(GS-4$'$) & 통과(verifybox \S\ref{sec:ch6_S1},\ref{sec:ch6_dae}) \\
Plan A(Fredholm)를 M1--M5$\cdot\varepsilon$ 가속 후보로 제한; broad-kernel/Volterra 한계 명시 & 통과(\S\ref{sec:ch6_closure}) \\
저속 OCV 골격(O1--O6); $w_j$ vs $\rho_G(g)$ 구분, $n^{\eff}w_j=RT/F$ 정합 & 통과(\S\ref{sec:ch6_ocvfit}) \\
GITT relaxation 을 $k_j$ 직접 측정으로 단정 X; $\chi$ vs $\eta_\ct$ 분리 & 통과(\S\ref{sec:ch6_gittfit}) \\
★ 다온도 Arrhenius 로 $\Delta H_a$ 기울기 식별; $\mathcal A_j$ 고정 조건 명시 & 통과(\S\ref{sec:ch6_arrhenius}) \\
$0.1\mathrm C/1.0\mathrm C$ 를 strict limit 아닌 proxy/stress test 로 & 통과(\S\ref{sec:ch6_crate}) \\
★ 도입 관찰 $\to$ 반증 N1--N4 데이터 시험으로 loop 폐합 & 통과(\S\ref{sec:ch6_falsify} loopbox) \\
calorimetry 없이 발열 파라미터 피팅 X(forward prediction) & 통과(\S\ref{sec:ch6_heat}) \\
피팅 편의(cap/clip/transform)를 본문서 배제, practicebox 격리 & 통과(\S\ref{sec:ch6_S6},\ref{sec:ch6_forbidden}) \\
충전/양극/full-cell 확장을 음극 방전 검증 이후로 보류 & 통과(\S\ref{sec:ch6_extension}) \\
모든 가정 AL-60$\sim$69 인용; 무근거 established 0; FLAGGED 정직 표기 & 통과(\S\ref{sec:ch6_al}) \\
``자명/clearly/쉽게/obviously'' 본문 0건(G-noleap) & 통과 \\
\bottomrule
\end{longtable}
}

% ====================================================================
\section{부록 B 가정 원장 (구 Chapter 6, Assumption Ledger, AL-60$\sim$69)}\label{sec:ch6_al}
% ====================================================================
본 장이 사용한 가정과 그 tier$\cdot$근거$\cdot$검증 DOI 다(통합 master 와 정합; AL-64$\sim$69 가 본 장 신규 등재).
{\small
\begin{longtable}{p{0.05\textwidth} p{0.49\textwidth} p{0.12\textwidth} p{0.26\textwidth}}
\toprule AL & 가정 진술 & tier & 근거(검증 DOI) \\ \midrule \endhead
60 & index-1 DAE(root-constrained ODE), BDF/Radau 시간 적분; 단조성 floor 가 index-1 보장 & GROUNDED & brenan1996(10.1137/\brk 1.9781611971224), hindmarsh2005(10.1145/\brk 1089014.1089020) \\
61 & Plan A closure: Fredholm 2종 ratio-substitution closed-form(사용자 PhD) & BOUNDED & jcp2017(10.1063/\brk 1.5000882), lee2011(10.1063/\brk 1.3565476), son2013(10.1063/\brk 1.4802584) \\
62 & ★ ratio ansatz 의 broad-kernel 열화 $=$ stretched-tail(저온) 영역서 부정확 $\to$ 직접 grid validator 필수 & BOUNDED & jcp2017 Eq.34 자기명시(10.1063/\brk 1.5000882) \\
63 & ICA($dQ/dV$) rate-dependency$\cdot$미분 잡음 증폭$\cdot$best-practice & GROUNDED & dubarry2022(10.3389/\brk fenrg.2022.1023555), fly2020(10.1016/\brk j.est.2020.101329) \\
64 & ★(신규) root-find/solver 종료 tolerance 는 물리 scale($\delta_q Q_\cell$, $\xi\in[0,1]$)$\cdot$기계 epsilon floor 로 정함(임의 상수 아님) & BOUNDED & brenan1996(10.1137/\brk 1.9781611971224); 수치해석 표준 \\
65 & ★(신규) 분포 적분 축약(adaptive quadrature/moment/Prony)은 기준 grid 대비 검증 시만; Prony $a_\ell\ge0,\tau_\ell>0$ 은 완전 단조성 보존 & BOUNDED & lindsey1980(10.1063/\brk 1.440530, KWW$\leftrightarrow$분포), johnston2006(10.1103/\brk PhysRevB.74.184430, 연속합) \\
66 & ★(신규) closure 검증 tolerance $\varepsilon_\tol=\max(\varepsilon_{\mathrm{disc}},\varepsilon_\noise)$; \emph{특정 숫자값은 데이터 의존 미확정} & FLAGGED(값) / BOUNDED(정의) & dubarry2022(10.3389/\brk fenrg.2022.1023555), fly2020(10.1016/\brk j.est.2020.101329); 값은 사용자 데이터서 산출 \\
67 & ★(신규) calorimetry 부재 시 Ch2/Ch4 발열 항은 forward prediction 으로만(전압 데이터에 발열 파라미터 피팅 X) & BOUNDED & bernardi1985(10.1149/\brk 1.2113792), thomasnewman2003(10.1149/\brk 1.1531194) \\
68 & ★(신규) 식별성: $\{R_n,k_j^{\eff},w_j,\rho_G,Q_\bg,Q_{j,\tot},k_{j,\lim}\}$ 동시 자유화 금지; 순차 제약(OCV$\to$GITT$\to$Arrhenius$\to$C-rate)으로만 식별 & BOUNDED & bardfaulkner(ISBN 978-0-\brk 471-04372-0), fly2020(10.1016/\brk j.est.2020.101329) \\
69 & ★(신규) Arrhenius 분해: $\mathcal A_j$ 고정 기울기 $=-(\Delta H_a-\chi_j\mathcal A_j)/R$ = \emph{유효} 엔탈피($\chi\mathcal A/RT$ 가 $1/T$-선형 기여); 순수 $\Delta H_a$ 는 $\mathcal A_j\to0$ 외삽/보정; $k_0,\Delta S_a$ 절편 공선형($\kappa$ 미지 잔여) & BOUNDED & eyring1935(10.1063/\brk 1.1749604), bardfaulkner(ISBN 978-0-\brk 471-04372-0) \\
\bottomrule
\end{longtable}
}
\textbf{tier 요약.} GROUNDED: AL-60,63,69(분해). BOUNDED: AL-61,62,64,65,67,68,69(분리). \textbf{FLAGGED: AL-66 의
$\varepsilon_\tol$ \emph{특정 숫자값}}(정의는 측정 기반으로 BOUNDED, 값만 데이터 확보 후 확정). 본 장은 무근거
established 수치 임계값을 본문에 두지 않는다(consolidated\_ch6 의 $\varepsilon_\tol$ 허위를 \S\ref{sec:ch6_epsilon}
에서 정정).

% ====================================================================
\section{부록 B 의 결론 (구 Chapter 6)}\label{sec:ch6_concl}
% ====================================================================
첫째, \textbf{본 부록(구 Chapter 6)은 유도식(Ch1--5)을 ICA/DVA 데이터에 \emph{피팅하는 방법론}이다}. 피팅 평가 순서 S1--S6
(\S\ref{sec:ch6_S})이 한 점의 $\dd Q/\dd V_n$ 을 주는 forward 계산이고, 이를 측정 ICA 와 비교해 파라미터를
식별한다. \textbf{수치 해법(전하 보존식 root-find, index-1 DAE 시간 적분, self-consistent Volterra 의 직접 grid
풀이)은 이 피팅의 inner-loop solver 로 \emph{종속}}되며 주역이 아니다(사용자 5-30 의도 반영).

둘째, \textbf{피팅의 본질은 회귀가 아니라 순차 식별}이다(\S\ref{sec:ch6_identif}). $R_n$ 과 $\chi_j$ 처럼 같은
꼬리를 만드는 공선형 파라미터는 한 데이터에서 분리되지 않으므로, 저속 OCV$\to$GITT$\to$다온도 Arrhenius$\to$
C-rate series 를 순서대로 걸어 식~\eqref{eq:ch6_no_free} 의 동시 자유화 금지 집합을 차례로 닫는다. $R_n\ne R_\ct
\ne R_{n,\eff}$ 의 세 저항 분리(Ch1$\cdot$3$\cdot$4 계승)가 식별성과 발열 예측을 동시에 지킨다.

셋째, \textbf{전하 보존식 root 의 dynamic/OCV 두 형태는 derivative 가 다르므로 분리}하며(식~\eqref{eq:ch6_dyn_deriv}
vs \eqref{eq:ch6_ocv_root}), 후자의 분모는 Ch.2 평형 chemical capacitance 와 정합한다. index-1 DAE 는 단조성
floor 가 index 를 1 로 만든다는 \emph{유도}로 정당화되며, 강구동 stiffness 는 cap 이 아니라 implicit solver 와
물리 병목으로 처리한다.

넷째, \textbf{closure 의 기준 해법은 직접 $g$-grid}(\S\ref{sec:ch6_grid})이고, adaptive quadrature/moment matching/
Prony-rational/Plan A(Fredholm)는 모두 기준 해 대비 검증될 때만 쓴다. Plan A 는 문헌적 가속 후보이되 흑연의 넓은
spectrum(저온 stretched tail)에서 가장 부정확할 수 있어(\AL{62}, jcp2017 자기명시) 항상 직접 grid 가 validator 다.

다섯째, \textbf{★ consolidated\_ch6 의 $\varepsilon_\tol$ 허위(근거 없는 수렴 임계값, flagbox 0건)를 정정}했다
(\S\ref{sec:ch6_epsilon}). $\varepsilon_\tol$ 을 ``작으면 좋다''가 아니라 기준 해의 이산화 오차와 측정 잡음으로
인한 ICA 불확실도의 \emph{큰 쪽}(식~\eqref{eq:ch6_epsilon})으로 \emph{측정 기반} 재정의하고, 그 \emph{특정 숫자값}은
데이터 의존이라 FLAGGED 로 정직 표기했다. root-find$\cdot$solver tolerance 도 같은 원리로 물리 scale$\cdot$잡음
floor 에서 근거화했다(\AL{64}, GS-4$'$).

여섯째, \textbf{반증(forward-only)이 모델의 가치}다(\S\ref{sec:ch6_falsify}). Arrhenius 직선성(N1), 전위 의존
spectrum shift(N2), 저온 lineshape(N3), $\rho_G$ 역산 금지(N4)로 도입 관찰을 \emph{기각 가능}하게 시험한다.
calorimetry 부재 시 발열 항은 피팅하지 않고 forward prediction 으로만 둔다(\AL{67}).

일곱째, \textbf{도입 관찰(저온$\cdot$고율 긴 겹침 꼬리)} $\to$ 인과 사슬(Ch1) $\to$ 발열/속도론/히스테리시스
(Ch2--5) $\to$ 피팅식과 순차 식별$\to$ 반증 가능 데이터 시험(본 부록 N1--N4)으로 \emph{한 바퀴가 닫힌다}. 음극
방전 모델이 먼저 식별$\cdot$검증된 뒤에야 충전 branch$\cdot$양극$\cdot$full-cell 로 확장한다.

\section*{Chapter 2--5 로의 전달}
본 장이 확정한 상태변수($V_n,\xi_j,k_j,\xi_{\eq,j}$)와 식 위에서: Ch.2 는 \emph{전지 가역 반응열}을 평형
전위 온도계수로 잇고, Ch.3 은 $k_j$ 를 forward/backward 반응속도론(Level B)으로 확대하며, Ch.4 는 그 위에서
발열을, Ch.5 는 충방전 히스테리시스를 다룬다. \emph{(기존 Ch.6 의 피팅 실무$\cdot$수치 내용은 본 장
\S\ref{sec:ch1_simplefit}(심플 근사식)$\cdot$\S\ref{sec:planB}$\cdot$\S\ref{sec:planA}(Plan B/A closure)$\cdot$
\S\ref{sec:ch1_numeric}(수치 평가) 및 \emph{부록 B}(\S\ref{sec:ch1_fitting_practice}: 식별성$\cdot$수치 해법$\cdot$순차 제약$\cdot$반증 워크플로)로 흡수했다 — Ch1 자기완결. Ch.6 은 해체되었다.)}

\begin{thebibliography}{99}
\bibitem{doyle1993} M. Doyle, T. F. Fuller, J. Newman, ``Modeling of Galvanostatic Charge and Discharge
  of the Lithium/Polymer/Insertion Cell,'' \emph{J. Electrochem. Soc.} \textbf{140}, 1526 (1993).
  DOI: 10.1149/1.2221597.
\bibitem{newman} J. Newman, K. E. Thomas-Alyea, \emph{Electrochemical Systems}, 3rd ed., Wiley (2004).
  ISBN 978-0-471-47756-3.
\bibitem{mckinnon1983} W. R. McKinnon, R. R. Haering, ``Physical Mechanisms of Intercalation,'' in
  \emph{Modern Aspects of Electrochemistry} No.~15, Plenum, New York (1983), p.~235.
\bibitem{bazant2013} M. Z. Bazant, ``Theory of Chemical Kinetics and Charge Transfer Based on
  Nonequilibrium Thermodynamics,'' \emph{Acc. Chem. Res.} \textbf{46}, 1144 (2013).
  DOI: 10.1021/ar300145c.
\bibitem{eyring1935} H. Eyring, ``The Activated Complex in Chemical Reactions,'' \emph{J. Chem. Phys.}
  \textbf{3}, 107 (1935). DOI: 10.1063/1.1749604.
\bibitem{evanspolanyi1938} M. G. Evans, M. Polanyi, ``Inertia and Driving Force of Chemical Reactions,''
  \emph{Trans. Faraday Soc.} \textbf{34}, 11 (1938). DOI: 10.1039/TF9383400011.
\bibitem{bardfaulkner} A. J. Bard, L. R. Faulkner, \emph{Electrochemical Methods}, 2nd ed., Wiley (2001).
  ISBN 978-0-471-04372-0.
\bibitem{marcus1956} R. A. Marcus, ``On the Theory of Oxidation-Reduction Reactions Involving Electron
  Transfer. I,'' \emph{J. Chem. Phys.} \textbf{24}, 966 (1956). DOI: 10.1063/1.1742723.
\bibitem{onsager1931} L. Onsager, ``Reciprocal Relations in Irreversible Processes. I,'' \emph{Phys. Rev.}
  \textbf{37}, 405 (1931). DOI: 10.1103/PhysRev.37.405.
\bibitem{degrootmazur1962} S. R. de Groot, P. Mazur, \emph{Non-Equilibrium Thermodynamics},
  North-Holland (1962); Dover (1984).
\bibitem{dahn1991} J. R. Dahn, ``Phase Diagram of Li$_x$C$_6$,'' \emph{Phys. Rev. B} \textbf{44}, 9170
  (1991). DOI: 10.1103/PhysRevB.44.9170.
\bibitem{ohzuku1993} T. Ohzuku, Y. Iwakoshi, K. Sawai, ``Formation of Lithium-Graphite Intercalation
  Compounds...,'' \emph{J. Electrochem. Soc.} \textbf{140}, 2490 (1993). DOI: 10.1149/1.2220849.
\bibitem{funabiki1999ea} A. Funabiki, M. Inaba, T. Abe, Z. Ogumi, ``Nucleation and Phase-Boundary
  Movement upon Stage Transformation in Lithium-Graphite Intercalation Compounds,''
  \emph{Electrochim. Acta} \textbf{45}, 865 (1999). DOI: 10.1016/S0013-4686(99)00290-X.
\bibitem{funabiki1999jes} A. Funabiki, M. Inaba, T. Abe, Z. Ogumi, ``Stage Transformation of
  Lithium-Graphite Intercalation Compounds Caused by Electrochemical Lithium Intercalation,''
  \emph{J. Electrochem. Soc.} \textbf{146}, 2443 (1999). DOI: 10.1149/1.1391953.
\bibitem{baessler1993} H. B\"assler, ``Charge Transport in Disordered Organic Photoconductors,''
  \emph{Phys. Status Solidi B} \textbf{175}, 15 (1993). DOI: 10.1002/pssb.2221750102.
\bibitem{svare2000} I. Svare, S. W. Martin, F. Borsa, ``Stretched Exponentials with $T$-Dependent
  Exponents from Fixed Distributions of Energy Barriers for Relaxation Times in Fast-Ion Conductors,''
  \emph{Phys. Rev. B} \textbf{61}, 228 (2000). DOI: 10.1103/PhysRevB.61.228.
\bibitem{johnston2006} D. C. Johnston, ``Stretched Exponential Relaxation Arising from a Continuous Sum
  of Exponential Decays,'' \emph{Phys. Rev. B} \textbf{74}, 184430 (2006). DOI: 10.1103/PhysRevB.74.184430.
\bibitem{lindsey1980} C. P. Lindsey, G. D. Patterson, ``Detailed Comparison of the Williams-Watts and
  Cole-Davidson Functions,'' \emph{J. Chem. Phys.} \textbf{73}, 3348 (1980). DOI: 10.1063/1.440530.
\bibitem{jcp2017} (사용자 PhD) K. Lee, S. Lee, C. H. Choi, S. Lee, ``Effects of External Electric Field and Anisotropic Long-Range Reactivity on Charge Separation Probability,''
  \emph{J. Chem. Phys.} \textbf{147}, 144111 (2017). DOI: 10.1063/1.5000882.
\bibitem{lee2011} (Refs.~6) S. Lee, C. Y. Son, J. Sung, S.-H. Chong, ``Communication: Propagator for Diffusive Dynamics of an Interacting Molecular Pair,''
  \emph{J. Chem. Phys.} \textbf{134}, 121102 (2011). DOI: 10.1063/1.3565476.
\bibitem{son2013} (Refs.~7) C. Y. Son, J. Kim, J.-H. Kim, J. S. Kim, S. Lee, ``An Accurate Expression for the Rates of Diffusion-Influenced Bimolecular Reactions with Long-Range Reactivity,''
  \emph{J. Chem. Phys.} \textbf{138}, 164123 (2013). DOI: 10.1063/1.4802584.
\bibitem{dubarry2022} M. Dubarry, D. Anseán, ``Best practices for incremental capacity analysis,''
  \emph{Front. Energy Res.} \textbf{10}, 1023555 (2022). DOI: 10.3389/fenrg.2022.1023555.
\bibitem{fly2020} A. Fly, R. Chen, ``Rate dependency of incremental capacity analysis (dQ/dV)...,''
  \emph{J. Energy Storage} \textbf{29}, 101329 (2020). DOI: 10.1016/j.est.2020.101329.
\bibitem{brenan1996} K. E. Brenan, S. L. Campbell, L. R. Petzold, \emph{Numerical Solution of Initial-Value
  Problems in Differential-Algebraic Equations}, SIAM Classics in Applied Mathematics 14 (1996).
  DOI: 10.1137/\brk 1.9781611971224.
\bibitem{hindmarsh2005} A. C. Hindmarsh, P. N. Brown, K. E. Grant, S. L. Lee, R. Serban, D. E. Shumaker,
  C. S. Woodward, ``SUNDIALS: Suite of Nonlinear and Differential/Algebraic Equation Solvers,''
  \emph{ACM Trans. Math. Softw.} \textbf{31}, 363 (2005). DOI: 10.1145/\brk 1089014.1089020.
\bibitem{bernardi1985} D. Bernardi, E. Pawlikowski, J. Newman, ``A General Energy Balance for Battery
  Systems,'' \emph{J. Electrochem. Soc.} \textbf{132}, 5 (1985). DOI: 10.1149/\brk 1.2113792.
\bibitem{thomasnewman2003} K. E. Thomas, J. Newman, ``Thermal Modeling of Porous Insertion Electrodes,''
  \emph{J. Electrochem. Soc.} \textbf{150}, A176 (2003). DOI: 10.1149/\brk 1.1531194.
\end{thebibliography}

\end{document}
